\chapter{CAMPO ALFA}\label{campo-alfa}

\hfill\break

Non c'era nessun ascensore spaziale su Campo Alfa, così facemmo il
viaggio verso la superficie su uno shuttle kristang. Solo che, ho
imparato, non si chiamava shuttle, navetta, o dropship, navicella da
sbarco, o lander, veicolo d'atterraggio, che erano altri termini che
avevo sentito usare dalla gente. Il termine utilizzato dall'esercito per
questi velivoli era Surface to Orbit Craft³ o Stoc. In pratica,
chiamavamo queste navi sia dropships, che Stocer, pronuncia ``stocher'',
come stock car, ma tutto attaccato. Le Stocer potevano raggiungere
l'orbita da sole senza bisogno di razzi ausiliari e alcuni modelli più
grandi erano in grado di effettuare voli interplanetari, se avevi un
sacco di tempo per le mani e andavi molto, molto d'accordo con le
persone con cui eri stipato là dentro. I kristang e altre specie avevano
velivoli che chiamavano qualcosa come dropship, ma le dropships venivano
utilizzate per invasioni su larga scala, per trasportare un sacco di
truppe e attrezzature fino alla superficie, un viaggio di sola andata.
Tecnicamente, una dropship era come un aliante della seconda guerra
mondiale, poteva sopravvivere all'ingresso in un'atmosfera, fare un po'
di manovre e atterrare in verticale, ma una volta a terra restava lì. Le
dropships erano veicoli irripetibili, mentre una Stocer poteva risalire
in orbita e fare il viaggio più volte. La distinzione non m'importava
molto, perché sarei stato solo un passeggero, ma era bello sapere come
facevano le cose i nostri soci anziani, mi faceva sentire parte della
squadra, non solo un ingranaggio in una macchina.

La navicella, o dropship, in cui m'infilai a forza conteneva
duecentoventi persone e la nostra attrezzatura, aveva più o meno le
dimensioni di un 747 e la maggior parte dello spazio era occupato da
motori e serbatoi di carburante. La discesa fu dura, colpimmo
l'atmosfera e mi sembrò di avere un sacco di cemento sul petto. Poi le
cose peggiorarono. Ragazzi, fui contento quando la nave si stabilizzò e
per un po' volò come un aereo. Quando successe applaudimmo tutti.
L'atterraggio vero e proprio fu tranquillo, capii che eravamo a terra
solo quando il rumore del motore s'interruppe.

Io. Che atterro su un pianeta alieno. Chi l'avrebbe mai detto, eh?
«Tenente», chiese qualcuno, «sa chi è stato il primo umano a mettere
piede su un altro pianeta?»

«Probabilmente il generale Meers», disse un soldato dietro di me. Meers
era al comando dell'Unef⁴.

«No.» Gonzalez scosse la testa. «Era un capitano donna dell'esercito
britannico, faceva parte della squadra di avanguardia multinazionale.
Era previsto che a uscire per primo fosse un colonnello dell'esercito
cinese, ma la capitana inglese era seduta vicino alla porta e, quando
atterrarono, il pilota kristang disse loro di uscire, subito. I kristang
non lo consideravano un avvenimento importante -- il primo atterraggio
dell'umanità su un altro pianeta -- volevano solo che fossero tutti
fuori dalla nave in fretta, in modo da poter tornare indietro e prendere
il gruppo successivo.»

Cazzo. Avevo sperato che la prima persona fosse stata un americano. Beh,
avevamo pur sempre Neal Armstrong, e il vecchio Neal ce l'aveva fatta
con una buona tecnologia americana. Questa britannica aveva solo chiesto
un passaggio su un bus spaziale alieno.

Il mio primo passo su un altro pianeta fu quasi imbarazzante, mi ero
abituato all'85\% di gravità della nave madre thuranin, poi di nuovo
all'assenza di gravità, e la gravità su Campo Alfa era al 107\% della
norma terrestre. Il peso inconsueto mi fece inciampare lungo la rampa.
Il tenente Gonzalez non scherzava sul caldo: quando scendemmo dalla
navicella, fui colpito da un'esplosione di calore che pareva uscire da
una fornace. Ed era afoso, lasciate perdere tutte le stronzate sul fatto
che il caldo secco si sentirebbe di meno. E luminoso, dovetti schermarmi
gli occhi e strizzarli per resistere, mentre percorrevo i cento metri
che ci separavano dal punto in cui erano parcheggiati i veicoli da
trasporto. Tutto odorava di bruciato, come in conseguenza di un
incendio. Il contenuto di ossigeno dell'atmosfera doveva essere solo un
po' più basso della norma terrestre, ma di certo avevo l'impressione che
i miei polmoni stessero succhiando aria e non ne stessero ottenendo
grandi benefici. Forse il punto di atterraggio era ad alta quota. Avevo
quasi raggiunto un camion da trasporto ed ero in fila per salire a
bordo, quando una donna dietro di me ondeggiò e rischiò di crollare sul
terreno polveroso. Anche a me la testa nuotava per il caldo improvviso,
come per un capogiro quando ti alzi troppo in fretta in spiaggia. Lei
protestò e cercò di allontanarmi, ma era evidente che le sue ginocchia
vacillavano. «Ascolta», lessi il nome sulla sua tasca sinistra, «ehi,
Miller, noi siamo Army Strong.» Ripetei uno slogan di reclutamento
dell'esercito americano: «Non lasciamo indietro nessuno, non in
combattimento e di certo non in questo caldo puzzolente, capito?».

Lei annuì con rabbia e si appoggiò a me, respirando a fatica, finché non
si alzò. Si sarebbe detto che avere bisogno di assistenza, anche per
poco, stesse uccidendo il suo orgoglio. «Grazie», guardò il mio
cartellino, «Bishop. Accidenti, fa caldo qui. Io sono di Seattle.»

«Ah! Qui Nord Maine. Non sai cos'è il freddo.» Lasciò che la tenessi per
le spalle per sorreggerla qualche secondo, finché le sue gambe non
smisero di tremare. A tre metri di distanza, le ginocchia di un tizio
cedettero di colpo e la sua faccia si piantò nella terra, un'altra
vittima del caldo improvviso. Miller si riprese e salì da sola sul retro
del camion. Una volta che il mezzo fu carico, ci dirigemmo verso una
tendopoli che stavano costruendo in cima a un promontorio. Mi guardai
attorno cercando di assaporare i miei primi momenti su un mondo alieno.
Non c'era molto da assaporare. Era asciutto e polveroso e l'aria odorava
ancora di bruciato, quindi la puzza non proveniva dallo scarico del
motore. Bassa vegetazione, raggruppata attorno alle rocce, sembrava
essere un tipo di lichene polposo e carnoso. Mi figurai che la pianta
immagazzinasse la propria riserva d'acqua, come un cactus. Rimbalzando
nel retro del camion, ci guardavamo l'un l'altro e ci scambiavamo
smorfie d'intesa. Campo Alfa era il nostro primo pianeta alieno. E
faceva schifo.

\hfill\break

Il tempo trascorso seduto sul camion fu l'unico che ebbi a disposizione
per godermi l'esperienza, non appena arrivammo al campeggio a noi
assegnato, l'esercito ci mise al lavoro per piantare le tende. Il caldo
non era così terribile, ora che avevo avuto il tempo di abituarmici, ma
la gravità in eccesso era stronza e la puzza di bruciato mi aveva
saturato il naso. Un sergente scelto assegnò me, Cornpone e altri cinque
all'allestimento di una grande struttura medica, fatta di pannelli
prefabbricati isolanti. Era più dura di quanto sembrasse. «Cominciamo
con questa parete laterale», dissi nella speranza che, completata la
prima struttura, saremmo potuti andare a mangiare.

Ci fu un suono ovattato, come una scoreggia molto lunga e acuta, poi
cominciò il gemito.

«È cornamusa? Che cazzo è?», esclamò Cornpone.

«Sì, c'è un battaglione inglese dall'altra parte del campo.» Alzai di
scatto il pollice al di sopra della mia spalla, indicando il punto in
cui avevo visto gli inglesi costruire il loro accampamento.

«E suonano la cornamusa? Davvero? Pensavo fosse una cosa irlandese.»

Garibaldi lo derise. «Scozzese, non irlandese, testa di rapa. Cazzo, ma
non sai niente?»

Valdez disse: «Ehi, hai sentito quella del tizio scozzese? Se ne sta
seduto in un bar ad affogare i dispiaceri nel whisky».

«Ce l'avete tutti?», chiesi. «Sollevate con le gambe, non con la
schiena.»

Cornpone sbuffò: «Accidenti, Bish, parli come mia madre. Lascia fare a
noi».

«Così dice al barista», Valdez grugnì per lo sforzo, «dice così:
``Costruisci trenta case, e quando cammini per la strada la gente dice:
`Ecco MacDougal, il costruttore di case'? No, non lo fanno''.»

Tenevamo la parete sulle ginocchia e io dissi alla squadra di andarci
sotto e caricarla sulle spalle, la gravità in eccesso mi stava
uccidendo.

«Poi dice: ``Tu salvi cinque bambini da un edificio in fiamme, e quando
cammini per la strada la gente dice: `Ecco MacDougal, il soccorritore'?
No, non lo fanno''.»

«Valdez, vuoi star zitto un minuto? Bene, gente, oh, issa!», dissi nella
miglior voce autorevole da esercito americano che mi veniva.

Ci eravamo messi quella stronza di parete sulle spalle e stavamo
vincendo la nostra battaglia contro la gravità, quando Valdez sussurrò:
«Poi MacDougal dice: ``Ma ti scopi una pecora\ldots''».

La perdemmo. Gli uomini si dispersero mentre la parete si schiantava e
tutti si rotolavano a terra dalle risate. Cercai di abbaiare un ordine,
ma ridevo così forte che mi uscì una bolla di moccio dal naso e, quando
gli altri lo videro, si piegarono di nuovo in due dalle risate.

Il tenente Gonzalez ci guardò male e un sergente cominciò a camminare
nella nostra direzione. Era ufficiale: non avevo alcun diritto di essere
a capo di nulla, l'esercito aveva fatto un errore enorme dandomi anche
solo la responsabilità di montare una tenda. Radunai gli uomini e
provammo di nuovo a sollevare la parete, ma ogni volta che l'avevamo
sollevata a pochi centimetri da terra, Valdez iniziava a ridere della
sua stessa battuta e faceva scoppiare a ridere tutti. Alla fine, mandai
Valdez a prenderci dell'acqua, chiamai più di una mezza dozzina di
uomini, raddrizzammo quella cazzo di cosa e la sistemammo.

«Unapecora.» Cornpone rise e non potei fare a meno d'imitarlo.

\hfill\break

Dopo aver sistemato la tenda, io e Cornpone entrammo finalmente in
contatto con gli altri due ragazzi del nostro vecchio gruppo di fuoco:
il Sergente Greg Koch e il soldato Dave Czajka. Chiamavamo Koch
``Sergente'', era un soprannome, ma ti servivano almeno i galloni da
sergente per usarlo, a meno che non volessi che un georgiano di colore,
alto un metro e ottanta e con la testa rasata, ti fulminasse con lo
sguardo. E tu non lo volevi. Dave lo chiamavamo ``Ski'' perché era
originario di Milwaukee, di Polonia per la precisione, e anche se il suo
cognome non finiva per ``-ski'', noi pensavamo che avrebbe dovuto.
Davvero, era stata sua la colpa di non avere un soprannome decente
quando entrò a far parte della nostra unità. Ragazzi, era bello
rimettere insieme la banda! A ciascuno di loro tre avrei affidato la mia
stessa vita. In Nigeria eravamo finiti nella merda, ma eravamo riusciti
a far sì di guadagnare solo qualche cicatrice e alcuni brutti ricordi da
quella situazione. Prima di andarmene di casa, ero rimasto in contatto
con loro tramite sporadiche e-mail, ma non li avevo più sentiti da
allora e Cornpone e io non avevamo la certezza che avessero intrapreso
il nostro stesso viaggio dalla Terra. Erano su un'altra nave kristang ed
erano arrivati su Campo Alfa poche ore prima di noi. Koch ci fece
sistemare in una tenda che avevamo dovuto aiutare a montare e capì dove
andare a mangiare. Eravamo dei campeggiatori felici.

Il Sergente Koch andò a fare importanti cose da sergente, e la sala
mensa non avrebbe aperto per cena prima di altre tre ore, così Cornpone
scavò nel suo zaino alla ricerca di snack.

«Che cos'hai di buono?», chiese Ski.

Cornpone fece l'occhiolino. «Ehm, vediamo. Cazzo! Ho un'ampia scelta di
genuine barrette energetiche Hooah!⁵.»

«Un'ampia scelta?», chiesi scettico. «Davvero?»

«Oh, sì, un'ampia scelta a dir poco. Potrebbe essere una pletora, forse
anche, oh, un'autentica, dannata cornucopia di snack.» Fu divertente
sentire Jesse pronunciare ``pletora'' nel suo accento del profondo Sud.

«Basta con le stronzate, amico, mi stai facendo venire fame», protestò
Ski. «Che gusti hai?»

«Bene, ho cartone alla cannella, ovviamente, anche uvetta e ghiaia, e il
mio preferito, segatura originale.»

Merda. Mi toccò uvetta e ghiaia. Le uvette erano più dure da masticare
della ghiaia.

\hfill\break

La mattina dopo, ci fecero marciare fino al poligono di tiro dopo
mangiato. Faceva freddo, il termometro fuori dalla nostra tenda diceva
che quasi si congelava. L'eccitazione era alta: tutti ci aspettavamo di
mettere finalmente le mani su alcune incredibili armi super-tecnologiche
dei kristang. E forse vestiti che ti scaldavano e raffreddavano in
automatico. Mi sentivo nudo senza fucile. Di fronte a noi c'era un
sergente artigliere del corpo dei Marine, con un taglio militare così
severo che sembrava di potersi sfregiare la mano toccandolo. Indossava
pantaloni e stivali da combattimento, ma sopra la maglietta aveva una
felpa rossa dei Marine.

«Sono il sergente artigliere Cragen, del corpo dei Marine degli Stati
Uniti.» Cragen fece una pausa, mentre le grida d'incitamento dei Marine
-- ``Urrà'' e ``Semper fidelis'' -- risuonavano dalla folla. «Sono
sollevato nel sentire che abbiamo tra noi alcuni fucilieri qualificati.
Visto che il resto di voi è dell'esercito, cercherò di parlare
lentamente in modo che possiate capire.» Gli altri Marine risero di
quell'osservazione.

«Ascoltate, gente!» Su un tavolo di fronte a lui c'era una scatola
grigia, toccò un pulsante sul lato della scatola che si aprì senza
rumore. Tutti sporsero il collo allo stremo, avremmo finalmente visto le
armi super-tecnologiche con cui avremmo ucciso i criceti. Quindi è
comprensibile che dalla folla si fosse levato un brontolio quando il
sergente Cragen tirò fuori dalla scatola quello che sembrava un fucile
M4 standard.

«Oh, fanculo!»

«Non ci posso credere!»

«Stronzate, ragazzi!»

«Un merdoso M4?»

«Whisky Tango Foxtrot⁶?!», che significava, ovviamente: ``Che cazzo
è?''.

«Ragazzi, è una cagata», sussurrò Ski a me e a Cornpone. Quest'ultimo si
limitò a sputare sul terreno bruciacchiato e polveroso.

Questi furono alcuni dei commenti più blandi. Cragen preferì lasciare
che le proteste morissero da sé, piuttosto che strigliarci per mancanza
di disciplina. «Gente, salutate quello che noi umani chiamiamo M4b1. I
kristang non si fidano a darci in mano qualcosa di più potente, e non ci
servirà per le missioni che ci saranno assegnate nell'immediato futuro.
L'unità M4 Bravo spara munizioni simili alle calibro 5,56 Mike Mike⁷,
due volte le due tre cui siete abituati. La differenza è questa», prese
un proiettile dal tavolo e indicò la punta. «Questa è una carica
esplosiva. I kristang ci assicurano che queste munizioni penetreranno il
giubbotto antiproiettile che i soldati ruhar indossano, sempre che
selezioniate l'opzione esplosiva, perché l'esplosivo è una carica cava.
Come il proiettile Heat che usano i nostri carrarmati. È anche possibile
scegliere di disattivare l'esplosivo, cosa che farete, a seconda delle
regole d'ingaggio in vigore al momento. Per selezionare la modalità
``punta esplosiva'' usate questo interruttore qui», e indicò un pulsante
incassato sul retro del manico, cioè in posizione comoda sia per i
soldati destrorsi che per quelli mancini. «Fate scorrere questa
copertura sul manico per esporre questo pulsante incassato. È incassato,
quindi non sarete tanto idioti da attivare l'opzione ``punta esplosiva''
per sbaglio. Dopo aver sparato nove colpi, cosa che si può fare in
singoli colpi standard o in tripla raffica, è necessario premere di
nuovo il pulsante per selezionare la modalità esplosiva.»

Cragen si dilungò sul M4 Bravo, ma nella maggior parte delle persone
l'eccitazione era sparita. Dov'erano le munizioni autoguidate, i laser,
i miglioramenti genetici per farci diventare super-soldati, gli occhi
bionici, le armature che ci avrebbero permesso di saltare nove metri in
alto e correre a centosessanta chilometri orari? Dov'era tutta la roba
figa che avevamo visto nei film di fantascienza? Questa era una
stronzata! I ruhar potevano colpirci dall'orbita con maser, cannoni a
rotaia e missili intelligenti, e tutto quello che avevamo erano fucili
dell'era del Vietnam con petardi per proiettili?

Una stronzata totale.

\hfill\break

Un paio d'ore dopo, avevo un M4 Bravo in mano al poligono di tiro. Era
proprio come il mio vecchio, fidato M4, che avevo lasciato a Fort Drum
quando ero andato a casa in licenza. E non sono mai tornato a Drum,
quindi il mio fucile potrebbe essere ancora lì. I kristang ci fornirono
un kit di conversione per i nostri normali M4 e le munizioni con punte
esplosive. In seguito, avremmo appreso che gran parte delle nostre
munizioni provenivano direttamente dalla Terra, senza raffinati
miglioramenti, perché i kristang avevano una fornitura limitata di
materiale esplosivo. Ogni squadra avrebbe avuto lastrine kristang pari
alla metà della carica normale, il resto sarebbero state banali
munizioni 5,56 mm. I nostri alleati pensavano che non avremmo avuto
bisogno di munizioni costose per le missioni che ci stavano assegnando.
Era un po' offensivo, qualsiasi duro lavoro in questa guerra sarebbe
stato fatto dagli adulti, noi soldati umani eravamo seduti al tavolo dei
bambini.

Il mio M4b1 arrivò con due lastrine di munizioni a punta esplosiva.
Prenotai il mio turno al poligono di tiro, esercitandomi con il pulsante
che consentiva di selezionare o escludere l'esplosivo. La modalità
esplosiva era impressionante, ci fu grande esultanza lungo la linea di
tiro quando vedemmo bersagli metallici triturati. I nostri proiettili
normali rimbalzavano sui bersagli, lasciando solo piccole ammaccature
lucide. Un vantaggio delle munizioni kristang consisteva nel fatto che
ogni cartuccia era esattamente uguale all'altra, fino al nano livello;
anche meglio di quelle che i nostri cecchini chiamavano ``munizioni da
competizione''. E le munizioni kristang sparavano in modo molto più
fresco e pulito, tanto che potevi ballare il rock`n'roll senza fondere o
sporcare la canna. Poi c'era ancora da smontare l'arma e rimontarla.
Bendati. E pulirla. Presi uno dei proiettili kristang da una lastrina e
lo esaminai da vicino. Non fosse stato per il colore diverso e la
superficie perfettamente liscia, senza freghi né graffi, sembrava la
munizione Nato standard che usavo da anni.

Stavamo andando in battaglia senza nessuno dei miglioramenti sperati,
costretti a chiedere un passaggio per ogni pianeta e senza che ci
fossero state affidate missioni importanti. Era come se i kristang
stessero solo assecondando il nostro desiderio di buttarci nella
mischia. Questo era il futuro.

Non era quello che mi sarei aspettato.

\hfill\break

Quel primo pomeriggio in cui ci trovavamo al poligono di tiro, un gruppo
di aerei passò in volo abbastanza basso e ravvicinato da permetterci di
vederli distintamente. Il gruppo girò intorno per ore, praticando
atterraggi, decolli e voli in formazione. Dico aerei invece di veicoli
spaziali, perché gli istruttori al poligono di tiro ci dissero che si
trattava di quelli: velivoli progettati per volare in un'atmosfera e non
in grado di raggiungere l'orbita. Quello che mi sorprendeva era il fatto
che fossero aerei ruhar, non kristang. Ovunque andassimo, i kristang
volevano che i nostri piloti fossero pronti a usare l'equipaggiamento
sottratto ai criceti, piuttosto che essere costretti a trasportare il
proprio equipaggiamento attraverso i corridoi aerei stellari. Ce n'erano
due tipi di base, noti con i nomi dati dagli umani: la ``Poiana'', che
era una nave da trasporto simile a un convertiplano V-22 Osprey, e una
sorta di cannoniera a due posti, simile a un Apache, che chiamammo
``Pollo''. Così come chiamavamo i ruhar criceti, davamo nomi
irrispettosi ai loro aerei. Erano entrambi velivoli a decollo e
atterraggio verticale come elicotteri o convertiplani, ma, invece delle
pale del rotore, avevano le gondole del motore a reazione all'estremità
delle loro corte ali. Tutti dovemmo ammettere che apparivano fighi. Ed
eravamo tutti orgogliosi di sapere che erano gli umani a pilotarli, e
che li avremmo pilotati in combattimento. Presto, speravamo.

\hfill\break

Il giorno dopo, ricevemmo alcuni degli altri giocattoli minimamente
aggiornati che i kristang ci permettevano di avere. Il giocattolo più
utile era una radio tattica personale, che trasmetteva utilizzando salti
di frequenza e raffiche criptate, e utilizzava rumore di disturbo bianco
ad ampio spettro in sottofondo, così la sua posizione risultava molto
difficile da individuare. Ricevemmo tutti una radio personale delle
dimensioni di un tipico smartphone terrestre, non fosse che era sottile
come una carta di credito, con un touch screen inscalfibile che non
potevi danneggiare, a meno di sparargli contro un proiettile a punta
esplosiva. Con la radio ci diedero auricolare e microfono da indossare
sotto i caschi. Le telecamere che avevamo sul casco trasmettevano video
attraverso le nuove radio, così capisquadra, plotone e compagnia
potevano vedere quello che vedevano i soldati. Potevamo anche togliere
le radio dalla cintura e usare lo schermo per selezionare le telecamere
della nostra squadra e vedere quella di chiunque altro. Inoltre, al
posto di una sola persona per unità designata come operatore radio e
costretta a portarsi in giro un pesante equipaggiamento, ogni soldato
ora aveva la capacità di comunicare con chiunque altro, in tutto il
pianeta. Il touch screen apriva un menu scritto in qualsiasi lingua
umana si parlasse. Inglese, spagnolo, francese, mandarino, hindi,
qualsiasi cosa, la radio capiva tutto. Tutti dovemmo ammettere che era
roba forte. Le persone che avevano una specializzazione professionale
militare relativa alle apparecchiature di comunicazione erano
disoccupate, dal momento che gli esseri umani non erano autorizzati ad
armeggiare con l'ingranaggio e, comunque, non lo capivano. Per quanto mi
riguarda, era tutta magia per me.

L'esercito voleva che chiamassimo il nuovo dispositivo RadTat, che sta
per Radio tattica. Alcuni burloni lo battezzarono ``zPhone'', perché Z
era oltre la I nell'alfabeto, almeno quanto lo zPhone era oltre
l'iPhone. Sono certo che l'esercito aveva qualche nome ufficiale come
Radio tattica, multiuso, fornita dai kristang, in conformità con lo
standard Mil bla bla bla. Vabbè. La caratteristica più figa delle RadTat
era che si potevano utilizzare come traduttori, che era una grande cosa,
perché nessuno di noi parlava una sola parola di ruhar. Io parlavo un
po' di spagnolo sgangherato e avevo imparato alcune frasi utili in hausa
e yoruba quando ero in Nigeria, ma, di fatto, non ero nemmeno molto
bravo a parlare americano. Facevo ancora confusione tra your e you're,
il che faceva impazzire mia madre quando le mandavo e-mail con errori
grammaticali. Facemmo tutti pratica con le nostre RadTat in modalità
traduttore: non dovevi fare altro che mettere l'auricolare in un
orecchio e parlare a voce lenta e chiara nel piccolo microfono a braccio
vicino alla bocca. Dovevi parlare a lungo nella tua lingua quando usavi
la tua RadTat per la prima volta, in modo che il traduttore potesse
capire il tuo particolare schema vocale. In qualche modo, quel coso era
in grado di capire la differenza tra il mio accento del Nord-Est e la
parlata strascicata del Sud di Cornpone. Il che era di per sé
impressionante. Tu parlavi nel microfono, facevi una pausa e le tue
parole uscivano dall'altoparlante della RadTat in qualunque lingua
avessi selezionato. Era perfino in grado di rilevare un'altra lingua
parlata vicino a te e tradurla in automatico nella tua. Il suono in
entrata e in uscita era più fluido di quanto mi aspettassi, anche se era
ancora chiaro che si trattava di una voce computerizzata. Quello che mi
colpì era che la RadTat gestiva piuttosto bene lo slang, anche
l'onnipresente slang militare. Quando non capiva qualcosa, emetteva un
ronzio, l'equivalente di ``ehm''. A un paio di chilometri da noi su
Campo Alfa c'era un battaglione francese, due loro squadre vennero a
farci visita un giorno e i francesi furono bombardati di richieste di
parlare in modo da poter provare, sia noi che loro, le nostre RadTat in
modalità traduzione. Fu fantastico, loro cianciavano in francese nei
microfoni e io li sentivo in inglese. Imparai, per prima cosa, che in
Francia il french toast non si chiama french toast, si chiama pain
perdu. Il che la dice lunga, se ci pensate: cioè, perché i francesi non
lo chiamano semplicemente toast? E allora noi che chiamavamo il nostro
formaggio ``formaggio americano''? Perché non solo ``formaggio''? No?
Forse sto andando fuori tema. Comunque, era strano sentirmi parlare e
sentire la RadTat ripetere le mie parole in francese, soprattutto perché
la RadTat di fatto usava un accento francese. Era così spassoso, per
entrambe le parti, che ci divertimmo un mondo a dire cose stupide e
sentire la RadTat tradurle. È probabile che quello non fosse il tipo di
addestramento a cui l'esercito intendeva sottoporci, ma sfruttammo al
massimo il traduttore e acquisimmo dimestichezza col suo uso.

Alcune persone riferirono di avere provato il traduttore con le truppe
cinesi e indiane in visita nelle loro zone della base, e aveva
funzionato alla grande. Per quanto riguarda il funzionamento con il
ruhar, facemmo pratica con un computer che parlava quella lingua e tutto
ciò che posso dire è che di certo speravo che quell'aggeggio diabolico
funzionasse bene sul campo come accadeva nella pratica.

\hfill\break

Un altro nuovo giocattolo attirò ulteriori gemiti quando ci fu mostrato.
Gunny Cragen era di nuovo in piedi su un palco di fronte a noi, questa
volta tirò fuori da una scatola un grosso tubo. «Salutate il missile
Fgm-148 Bravo Javelin.»

Merda. Lo Javelin era un'arma abbastanza decente, un missile anticarro
spalleggiabile, lancia e dimentica, che avevo usato in Nigeria, ma aveva
dei problemi. Tanto per cominciare, la dicitura ``spalleggiabile'' era a
dir poco un eufemismo, l'ordigno infernale pesava ventidue chili. I
civili potrebbero pensare che un'arma che consente a un singolo soldato
di distruggere un carro armato principale valga ognuno dei suoi ventidue
chili, ma non ne hanno mai trasportato uno attraverso la boscaglia
giorno dopo giorno. E sono ventidue chili cui vanno aggiunti il fucile,
centottanta cartucce di munizioni, attrezzatura radio, borraccia, cibo
pronto e tutte le altre cose che l'esercito pensa tu debba portare.
Perciò, gli Javelin sono assegnati a squadre di due uomini, come minimo.
In battaglia, un uomo fa da osservatore, mentre l'altro aziona l'arma,
poi entrambi escono da lì il più in fretta possibile. ``Spara e
scappa'', lo chiamiamo noi.

Comunque, Cragen stava parlando: «Conoscete tutti i punti di forza e di
debolezza del missile Javelin. I nostri alleati ci hanno aiutato con il
principale punto debole: il sistema di guida elettro-ottico a
infrarossi. Il manuale dice che la guida a infrarossi si raffredda in un
paio di secondi, ma sappiamo tutti che dipende dalle condizioni in cui
ci si trova. Quando ci sono novantotto gradi nel deserto o nella
giungla, impiega più tempo ad attivarsi. I kristang ci hanno dato dei
kit per aggiornare la guida sul Fgm-148 Bravo: ora si raffredda del
tutto in meno di mezzo secondo e rimane fresco fino a mezz'ora, quindi
non dovete preoccuparvi se avete necessità di usarlo subito. Dopo
mezz'ora dovete spegnerlo, ma dieci minuti dopo potete attivarlo di
nuovo. Tenete presente che, ogni volta che lo utilizzate, non rimarrà
fresco a lungo, è tutto nel manuale. Anche il sistema di imaging della
guida è aggiornato, quindi se mai avete avuto problemi, in una giornata
calda, a far sì che l'arma riconoscesse un obiettivo caldo quasi come
ciò che lo circondava, questo non sarà più un problema. Inoltre, una
volta che abbiate identificato un bersaglio per l'arma, non lo perderà,
anche se vi muovete o si muove il bersaglio».

Ci fu una manciata di lenti, ironici applausi. Cragen ci gettò uno
sguardo aspro. «So che vi aspettavate fucili laser e bombe atomiche
grandi come una bomba a mano quando vi siete diretti verso le stelle.
Non ne abbiamo, questo è quello che abbiamo», si sollevò lo Javelin
Bravo sopra la testa. «E combatteremo al meglio delle nostre capacità.»

Quello che non disse è che la nostra capacità di combattere sarebbe
stata maggiore se avessimo avuto armi migliori. Vabbè. L'esercito dice
che combattiamo con la forza che abbiamo, non con la forza che vorremmo
avere. Hooah!

\hfill\break

Dopo una settimana di addestramento alla guerra spaziale un po'
deprimente e deludente, con, in sostanza, le stesse armi che avevo usato
in Nigeria, ricevemmo finalmente un nuovo giocattolo ad alta tecnologia
con cui trastullarci. Il missile anticarro Javelin Bravo, come il nostro
Javelin originale, aveva una capacità antiaerea minima, limitata all'uso
contro elicotteri in volo a bassa quota. I kristang dovevano aver deciso
che nessuno dei nostri missili da difesa aerea portatile, i Manpad⁸, era
utile in questo conflitto. Avevo usato i nostri missili Stinger in
addestramento, mai in combattimento. Lo Stinger era un missile
maledettamente efficace contro elicotteri, droni, aerei ad ala fissa e
persino jet a volte. Solo se eri fortunato e il jet non era andato fuori
portata prima che tu avessi preparato e puntato lo Stinger. Contro le
navi supersoniche che salivano in verticale in orbita, non vedevo come
uno Stinger avrebbe potuto aiutarci. Avevo visto persino i nostri Polli
catturati, ossia le cannoniere ruhar, passare a velocità incredibile dal
volo a punto fisso alla salita rapida in verticale. Tutti quelli che
conoscevo erano d'accordo che noi della fanteria dovevamo fare qualcosa
di serio contro la minaccia aerea, o qualsiasi situazione di
combattimento sarebbe finita prima di iniziare.

Il missile Manpad che i kristang ci fornirono lo chiamammo subito
Zinger, ovvio. Pesava circa tredici chili, non era molto diverso da uno
Stinger, la differenza è che il dispositivo di puntamento Zinger era
integrato nel tubo di lancio ed era monouso, mentre con il nostro
Stinger dovevi attaccare la scatola del dispositivo di mira al tubo
quando volevi lanciare il missile. Uno Zinger poteva salire fino a
ventiquattromila metri, era ipersonico, e se mancava un bersaglio poteva
guardarsi intorno e decidere d'ingaggiarne un altro da solo, oppure
l'utente poteva reindirizzare il suo dispositivo di mira tramite zPhone,
o il missile poteva comunicare con altri Zinger nella zona per ricevere
indicazioni sul bersaglio da puntare. Le sue ali potevano anche
dispiegarsi in modo che il missile si aggirasse nell'area fino a venti
minuti prima che il carburante fosse esaurito, quella caratteristica
significava che potevamo spararlo in anticipo in uno spazio aereo prima
che il nemico arrivasse, questo dava all'utente il tempo di andarsene.

Lo Zinger ricevette un'accoglienza genuina ed entusiastica quando ci fu
mostrato, e lo usammo davvero contro i droni. A differenza di quanto
accadeva nella pratica con gli Stinger, quando l'esercito statunitense
ci faceva sparare piccoli razzi invece di costosi missili, i kristang
non si preoccupavano che consumassimo le munizioni, anche se nei nostri
giri di prova le testate erano disattivate. Era un'immensa
soddisfazione, per noi soldati semplici, vedere i droni presi di mira
cadere dal cielo.

Poi divenne più dura. I droni cui sparammo all'inizio erano stupidi,
erano in grado di fare un po' di manovre per evitare i missili, ma
avevano capacità difensive minime. In seguito, facemmo pratica contro i
droni che simulavano aerei ruhar e navicelle come Polli, Poiane e il
brutto velivolo d'assalto, la navicella che chiamavamo ``Avvoltoio''.
Tutti avevano torrette difensive a laser o a fascio di particelle che
potevano confondere, bruciare o distruggere uno Zinger, anche se si
avvicinava a scatti per evitare le difese del fascio. Gli Zinger
dovevano avere a che fare anche con le contromisure elettroniche, e
queste non si limitavano allo spettro di infrarossi o microonde. Gli
aerei ruhar avevano una capacità di stealth⁹ attivo che rendeva
difficile avvistarli, anche all'interno dello spettro del visibile. Non
è che con il sistema stealth in funzione l'aereo diventasse davvero
invisibile, ciò che avveniva è che l'aria attorno al velivolo
s'increspava, distorcendone le forme, e luci multicolori si accendevano
e spegnevano a intermittenza in un involucro attorno al velivolo,
rendendo difficile stabilire con esattezza dove si trovasse. La tattica
che ci insegnarono a usare i kristang fu quella di sparare più Zinger
contro un solo bersaglio. In pratica, eravamo in grado di prendere un
lanciatore Zinger, aprire lo schermo di mira sul lato del tubo,
designare un bersaglio e fare fuoco in meno di un minuto. Un team Zinger
di due persone poteva trasportare e lanciare quattro Zinger, in teoria,
con una che azionava il sistema del missile e l'altra che faceva da
rilevatore. Non dimenticate che due Zinger pesavano più di ventisette
chili e, mentre trasportavi uno di quelli, si presume che fossi già
carico di fucile, munizioni, borraccia, zPhone, giubbotto antiproiettile
e qualunque altra cosa l'Unef si aspettava che un soldato vestito di
tutto punto si trascinasse dietro. Certo, mentre impostavi il secondo
colpo, l'aviazione nemica sarebbe stata sulle tue tracce e avrebbe
risposto al fuoco, cosa che tendeva ad accorciare l'aspettativa di vita
dell'utente e perciò l'efficacia bellica. L'ultimo commento non è una
battuta che mi sono inventato, per inciso, era scritto proprio nel
manuale d'istruzioni che ricevemmo dai nostri patroni, i kristang.
Avevano bisogno di lavorare sulle loro pubbliche relazioni, alla grande.

Comunque, con gli Zinger, noi soldati semplici dell'Unef sentimmo che
eravamo equipaggiati per la battaglia in modo quasi adeguato e non ci
saremmo dovuti limitare ad accovacciarci nelle trincee e morire in
combattimento. Potevamo rispondere al fuoco. Rispondere solo fino a
ventiquattromila metri, è vero, perché contro le navi in orbita non
potevamo fare niente. Ma ci dava speranza. Se i kristang ci avevano
affidato i loro giocattoli, e ce ne avevano forniti in abbondanza,
stavano pianificando di farci fare qualcosa di utile in combattimento.
Ovunque sarebbe stato.

\hfill\break

La maggior parte del nostro tempo di addestramento su Campo Alfa non era
spesa per imparare a usare le armi nuove o modificate. Era un
addestramento di aggiornamento per soldati come me, che erano dediti
alla raccolta delle colture o all'aiuto dei civili quando i ruhar
avevano attaccato. Il lavoro che avevo fatto con la guardia nazionale
nel Maine era stato duro, da spaccarsi la schiena, ma non mi aveva
preparato al combattimento. Prima di arrivare su Campo Alfa, l'ultima
volta che avevo sparato con un fucile, a parte a caccia, era stato in
Nigeria. Ero arrugginito e un po' fuori forma. Le escursioni con gli
zaini da quarantacinque chili, l'addestramento corpo a corpo e la corsa
a ostacoli mi rafforzarono, e ne fui felice. I ruhar non ci sarebbero
andati piano con me, era meglio essere dolorante ed esausto ora, che non
essere preparato quando saremmo entrati in azione.

Correre con quel caldo mi faceva a pezzi. Correre con la gravità in
eccesso mi faceva a pezzi. Mettete insieme le cose, e mi tremavano le
gambe dopo pochi chilometri. Cornpone cadde in ginocchio e vomitò le
budella tre volte durante la nostra prima corsa di otto chilometri. Dopo
di che, il Sergente Koch decise che ci saremmo alzati in anticipo, per
correre prima che facesse caldo. Su Campo Alfa, sembrava esserci un
interruttore on/off per la temperatura: quando la stella arrivava
all'orizzonte al mattino, faceva subito caldo. Al tramonto, la
temperatura crollava. La terza mattina ci vide scivolare fuori dalle
nostre cuccette per correre nel freddo antelucano. Una giornata su Campo
Alfa durava un po' più di ventotto ore, il nostro programma lasciava il
tempo per sette ore buone di griglia, anche alzandoci presto.

Senza il caldo diurno, dovevamo vedercela solo con la gravità in
eccesso. E la sabbia che si sollevava e puzzava di bruciato. E la
schiacciante delusione che uscire allo scoperto non significava essere
stati trasformati in super-soldati. Saremmo diventati super-soldati alla
vecchia maniera, annunciò il Sergente Koch, attraverso il duro lavoro.
Quelle parole erano fonte d'ispirazione fino quasi al segno dei tre
chilometri, quando annaspavamo in cerca di fiato.

Il nostro gruppo di fuoco non era l'unica squadra ad alzarsi presto per
l'esercitazione: la seconda mattina stavamo correndo, eravamo a circa un
chilometro e mezzo dalla base sulla via del ritorno, quando girammo
attorno al lato di una collina e, sull'altro versante, trovammo una
mezza dozzina di persone, anche loro di ritorno alla base. Senza che
nessuno avesse bisogno di dire nulla, accelerammo. E loro fecero la
stessa cosa. E l'altro gruppo alzò la posta. Presto facemmo uno sprint.
Andai in testa. Corro fin da quando ero bambino e mia madre mi portava a
correre con lei. L'altro gruppo aveva due velocisti, poi uno. Davanti a
noi, i sentieri convergevano, raggiunsi l'incrocio poco prima dell'altro
corridore, che poi superò anche me, e corremmo più forte che potevamo,
mentre le nostre squadre dietro di noi applaudivano.

Solo che non era un lui, era una lei: avrei dovuto capirlo dalla corta
coda di cavallo che le ballonzolava dietro la testa, se ci avessi fatto
attenzione. Cazzo. In un certo senso la conoscevo, era nel nostro
battaglione di supporto alla brigata, si chiamava Shauna, o qualcosa del
genere, ma non la vedevo da tipo un mese prima di tornare dalla Nigeria.
Merda, non sapevo che fosse così veloce. Si spinse in vantaggio, solo
trenta, sessanta centimetri, ma non riuscivo a raggiungerla, per quanti
sforzi facessi. Anche lei stava facendo fatica, aveva l'andatura
traballante e il respiro spezzato come il mio. Sfrecciammo oltre i due
grandi massi che segnavano il confine non ufficiale della base e io
crollai sulle ginocchia, e fui costretto a usare le mani per spingermi
di nuovo in piedi e superarlo camminando. Lei stava in piedi accanto a
me, con le mani sulle ginocchia, ansante. «Cazzo.»

«Merda sì'', risposi. «Stai cercando di uccidermi?»

«Ehi, eri tu a spingere me.»

Le tesi una mano. «Joe Bishop.»

«Sì, ti conosco, sei quello di Barney, giusto? Shauna Jarrett.» Ci
stringemmo la mano.

«Vuoi il mio autografo?» Stavo scherzando solo a metà. Tutta la faccenda
di Barney mi stava davvero stufando.

Lei inclinò la testa: «Perché, la faccenda di Barney ti ha fatto scopare
un sacco?».

Mi si accese una lampadina in testa. Sì, lo so, ogni volta che una
ragazza sorride a un ragazzo, lui spera che sia interessata, ma allora
pensai che lei potesse esserlo davvero. «Non così tanto», dissi
facendole l'occhiolino.

Lei rise, e fu bellissimo. «Non perdere la speranza, Joe.» Gli altri
componenti dei nostri gruppi arrivarono e diedero pacche sulla schiena a
Shauna per aver vinto quella gara improvvisata. Shauna e il suo gruppo
iniziarono ad andarsene.

Se ci fosse stata una qualche opportunità, non me la sarei lasciata
scappare. Eravamo molto, molto lontani dalla Terra, in guerra, e non
c'erano tutte queste donne nel giro di mille anni luce da Campo Alfa.
«Ehi, Shauna, ci vediamo in giro?»

Per la mia gioia, mi sorrise da sopra la spalla. «Forse?»

\hfill\break

Oltre a familiarizzare con il nuovo equipaggiamento e a rimetterci in
forma, trascorrevamo del tempo in classe, che significava sederci sul
pavimento sporco della tenda ad ascoltare lezioni sul nemico, sulla
struttura di comando dei nostri alleati e, cosa fondamentale, Le Regole.
Tutto mi sconvolgeva quando lo scoprivo per la prima volta, così lo
spiegherò lentamente. Farete meglio a mettervi comodi, perché ci sarà un
quiz a sorpresa alla fine.

Aspetterò mentre vi andate a prendere una birra.

Siete comodi?

Ok, prima di tutto, qualche informazione di base per voi. I kristang non
erano al vertice della coalizione alleata, cosa che avevo già immaginato
visto che dovevano fare un giro interstellare su una nave madre
thuranin. Quello che mi giungeva nuovo era che neanche i thuranin erano
al vertice. La coalizione alleata era guidata da alieni simili a gatti
chiamati maxolhx. Che fossero simili a gatti era una diceria, dato che
nessun umano ne aveva mai visto uno, e pochissimi kristang avevano avuto
quel privilegio. Poiché è probabile che un umano non vedrà mai un
maxolhx in vita sua, dirò solo che sono una specie antica, con una
tecnologia incredibile. Di certo, quando i maxolhx combattevano, erano
armati di qualcosa di meglio di un cazzo di M4. E so che maxolhx è un
nome strano, ma era la cosa migliore venuta in mente alla nostra gente
della G2 per trascrivere il suono della parola quando i kristang la
pronunciavano.

Il pianeta natale dei maxolhx si trovava a due terzi del percorso
intorno alla galassia dalla Terra, se questo vi dice qualcosa, che non
dice a me. È molto, molto lontano. Scoprimmo che i thuranin non erano
unici nel loro genere, erano responsabili solo di una piccola parte
della galassia sotto il controllo dei maxolhx e i kristang gestivano un
paio di migliaia di anni luce cubici sotto il controllo dei thuranin. I
maxolhx avevano altre specie di secondo livello oltre ai thuranin e i
thuranin avevano sotto il loro comando altre due specie, che noi
sapessimo, oltre ai kristang. I thuranin sono una sorta di omini verdi,
più bassi della media umana, ma umanoidi, con due gambe, due braccia,
due occhi e una testa calva. Gli esseri umani che avevano parlato con i
kristang avevano avuto l'impressione che a questi ultimi non piacessero
i loro mecenati e che il sentimento fosse reciproco: i thuranin
disprezzavano qualsiasi specie dotata di una tecnologia inferiore. Mi
chiedo cosa pensassero dei penosi umani.

Dalla parte del nemico, la situazione era simile. In cima c'erano i
rindhalu, che sono degli esseri dall'aspetto inquietante simili a ragni.
Il solo pensiero mi faceva accapponare la pelle. I loro alleati di
secondo livello, in questo quadrante della galassia, erano gli jeraptha,
che erano quasi altrettanto inquietanti, essendo insetti. Gli jeraptha
non sono una mente alveare come le api, sono più simili ai coleotteri. I
ruhar sono una specie cliente degli jeraptha, come i kristang sono una
specie cliente dei thuranin. Sulla base di alcune cose che dissero i
kristang, i nostri informatori della G2 ipotizzarono che i ruhar
avessero una tecnologia leggermente migliore di quella kristang.

La propaganda kristang, dico propaganda perché è quello che mi sembrava,
affermava che i rindhalu erano più vecchi, molto più vecchi, dei maxolhx
e avevano cercato di sopprimere lo sviluppo di altre specie intelligenti
nella galassia, finché i virtuosi maxolhx non si erano ribellati, molte
migliaia di anni prima. Gli alleati stavano combattendo per il diritto
delle specie intelligenti di evolversi da sole, mentre le specie nemiche
erano schiave del ragno inquietante rindhalu che voleva il controllo
totale della galassia. Per essere una specie avanzata, i kristang
avevano una propaganda molto goffa. Pensate all'Unione Sovietica o alla
Corea del Nord o alla Germania nazista. Noi alleati stiamo combattendo
il nemico malvagio per la gente comune gloriosa e virtuosa! Sotto la
magnifica guida dei nostri capi supremi, la nostra giusta causa
trionferà! L'unione fa la forza! E così via. Comunque, questo è quello
che ci dissero. Sono certo che la verità sia più complicata, lo è
sempre.

L'incredibile tecnologia dei maxolhx e dei rindhalu non si spingeva fino
a creare, o anche solo far funzionare, i wormhole che rendevano pratico
il volo interstellare a lungo raggio. Quei wormhole erano stati creati
milioni di anni prima, da una specie allora sconosciuta, se non per
quello che si era lasciata alle spalle. I kristang chiamavano i
costruttori di wormhole spaziali ``gli Anziani''. Ai wormhole non
importava chi li attraversasse, e non c'era modo conosciuto per
controllarne uno o spegnerlo. O danneggiarlo, neppure una testata
nucleare avrebbe avuto effetto. I wormhole funzionavano, ed è tutto
quello che ci serviva sapere.

Poi arrivammo alle Regole.

Le Regole erano le regole di ingaggio per il combattimento. Si
applicavano a entrambe le parti, venivano fatte rispettare con rigore
dai maxolhx e dai rindhalu. L'esistenza delle Regole spiegava come mai
la guerra andasse avanti da migliaia di anni, come mai entrambe le parti
non si fossero spazzate via a vicenda molto tempo prima. La base delle
Regole era che maxolhx e rindhalu non potevano permettersi di combattere
sul serio tra loro in modo diretto, come l'America e la Cina non
potevano combattersi a vicenda. Sulla Terra, potenze con megatoni di
testate nucleari non potrebbero scontrarsi senza distruggere i propri
paesi. Anche nella galassia, le specie dotate di armi che facevano
sembrare petardi le testate nucleari non potevano permettersi di farsi
la guerra in modo diretto. Perciò, combattevano attraverso le loro
specie clienti, che a loro volta avevano specie clienti, e così via. La
maggior parte dei combattimenti e dei morti era provocata tra specie di
terzo livello come i ruhar e i kristang. E tra loro clienti. Come gli
esseri umani.

Comunque, ecco Le Regole.

Farete meglio a prendere carta e penna.

Regola Numero 1: niente armi nucleari o altri tipi di armi radiologiche
su un pianeta abitato o in potenza abitabile o nelle sue vicinanze.
Neppure antimateria, perché un'arma del genere produce radiazioni forti,
anche se di breve durata. Questa regola era la prima della lista perché
i maxolhx e i rindhalu non volevano che i loro sudici subalterni
contaminassero un pianeta che avrebbero potuto un giorno desiderare per
loro stessi. Anche nella vastità della Via Lattea, il numero di pianeti
abitabili nel range pratico di un wormhole era limitato, quindi entrambe
le parti non volevano che pianeti utili fossero rimossi dall'uso
produttivo perché qualche idiota si beveva il cervello con le testate
nucleari.

Regola Numero 2: niente armi chimiche, per lo stesso motivo della Regola
Numero 1. Sostanze chimiche nocive potrebbero indugiare nell'ambiente
per molto, molto tempo. Non dicevano che ogni specie che abbandona un
pianeta è tenuta a spazzare alle proprie spalle e a portarsi dietro i
rifiuti, pensai che fosse sottinteso.

Regola Numero 3: niente nano-armi. La nano-tecnologia era utilizzata per
fare tante cose, ma non poteva essere utilizzata come arma. Persino i
kristang non erano sicuri della base di questa regola, ma si diceva che
ci fosse stato un incidente in cui la nano-tecnologia si era liberata e
aveva fatto sì che l'intero sistema stellare fosse messo in quarantena,
per sempre. Da allora, le nano-armi erano un enorme November Golf¹⁰ in
guerra, un disastro.

Regola Numero 4: niente armi biologiche. Questa regola non serviva tanto
a proteggere i maxolhx e i rindhalu, perché le specie erano così diverse
che un agente patogeno mortale per una non aveva per lo più effetto
sulle altre. Si trattava più di una sorta di Convenzione di Ginevra in
cui i combattenti accettavano di non fare qualcosa al nemico, in modo
che il nemico non lo facesse a loro. Altrimenti, ogni parte avrebbe
potuto semplicemente sganciare missili stealth con contenitori di armi
biologiche nell'atmosfera dei pianeti nemici e presto non ci sarebbe
stato nessun superstite da entrambe le parti.

La Regola Numero 4 risultava, ci dissero, un po' infida nell'uso reale,
perché virus, batteri e altri pericoli biologici mutavano nel corso del
tempo. E, se un virus naturale noto fosse diventato più contagioso o
mortale, chi avrebbe potuto dire che non fosse successo in modo
naturale? La Regola Numero 4 era in vigore per prevenire attacchi
biologici catastrofici, ma, a quanto pare, in una certa misura era
permesso imbrogliare, finché la situazione non sfuggiva di mano e
qualcuno non ci andava di mezzo. I kristang ci dissero che i ruhar erano
noti per aver oltrepassato i limiti della Regola Numero 4. Sono certo
che i ruhar avrebbero detto la stessa cosa dei kristang. La faccenda
delle armi biologiche mi spaventò, sebbene i kristang ci avessero
rassicurato che i ruhar non ne sapevano abbastanza di biologia umana per
far sì che un'arma divenisse efficace contro di noi. Quello che mi
spaventava era che i kristang avevano una tecnologia medica molto più
avanzata di quella che l'umanità aveva portato tra le stelle e non ne
condividevano nulla con noi. Ci dissero che, senza anni di studio della
biologia umana, non avrebbero potuto applicare in sicurezza la loro
tecnologia avanzata su di noi. Il che significava che anche un comune
virus influenzale, che qualche umano avesse portato con sé dalla Terra,
avrebbe potuto rappresentare un serio problema per la spedizione. E che
le truppe colpite avrebbero dovuto essere trattate con qualsiasi cura
medica disponibile su quel pianeta. Non ci sarebbe stato nessun volo di
evacuazione medica per gli Stati Uniti in quella campagna.

La Regola Numero 5 era, in poche parole, ``non lasciar cadere rocce''.
Niente asteroidi o comete indirizzate in modo da colpire un pianeta
abitabile. Questo significava anche non prendere nemmeno una piccola
roccia, accelerarla a velocità relativistica e sbatterla su un pianeta.
Ogni pianeta con una popolazione decente aveva sistemi in grado di
rilevare e deviare rocce che per natura si trovavano comunque su una
rotta d'impatto. La Regola serviva a impedire a un combattente di
saturare quelle difese o di racchiudere una roccia ad alta velocità in
un campo stealth. Che dire, potreste chiedere, dei cannoni a rotaia? Le
astronavi si servivano di norma di cannoni a rotaia per accelerare i
proiettili cinetici fino a velocità relativistiche e i cannoni a rotaia
erano usati sia per i combattimenti spaziali che per gli attacchi
planetari. La chiave nell'uso dei cannoni a rotaia per il bombardamento
planetario era l'effetto che avevano sul pianeta, bisognava limitare la
resa energetica di un singolo dispositivo cinetico a meno di un
megatone, e non si potevano usare così tanti colpi di cannone da
influenzare il clima del pianeta. Colpire un pianeta con un sacco di
cannoni a rotaia avrebbe potuto gettare una spessa nube di polvere in
aria e portare a un rapido raffreddamento; a quella polvere ci sarebbe
potuto volere molto tempo per stabilirsi fuori dall'atmosfera, anche
portando a una mini-era glaciale.

Quindi, queste erano Le Regole. L'idea di usare le maiuscole per queste
parole non è farina del mio sacco, c'erano sulle diapositive PowerPoint
dell'Unef. Il punto delle Regole sembrava essere questo: i pianeti
abitabili entro il range pratico di un wormhole sono rari e preziosi, e
maxolhx e rindhalu non potevano permettere a specie minori di
danneggiare beni immobili di valore. Le Regole non erano incentrate
sulla protezione di soldati o civili da entrambe le parti, non c'era
nessun equivalente alieno della Convenzione di Ginevra per la guerra
spaziale. Maxolhx e rindhalu mettevano in atto Le Regole per impedire
che la loro antica guerra sfuggisse loro di mano. Finché le specie
minori si conformavano alle Regole, potevano uccidersi a vicenda quanto
volevano.

A me tutte Le Regole andavano bene, specie perché noi dell'Unef non
avevamo nano-armi, né armi nucleari, biologiche, chimiche o
relativistiche con noi su Campo Alfa. Le Regole implicavano che queste
armi non sarebbero state usate contro di noi. Tranne forse quelle
biologiche, se un nemico pensava di poterla fare franca. Noi umani non
avevamo alcun motivo di preoccuparci per questo.

\hfill\break

Non vidi più Shauna correre al mattino, e non volevo rendermi sospetto
con uno squillo di zPhone. Fui entusiasta quando ci incontrammo per caso
in sala mensa, è così che l'esercito chiamava il refettorio. Era una
grande tenda, con teli sul terreno per tenere sporco e polvere lontani
dal cibo e ridurre al minimo l'odore acre di bruciato. Avevo un vassoio
con sandwich, burro d'arachidi e una specie di carne misteriosa con purè
di patate e quello che doveva essere sugo. E fagiolini. Veri fagiolini,
così sembrava, non erano stracotti. E un panino. Con il burro. Dopo
essermi fatto il culo tutto il giorno avevo fame, e tutto il cibo
sembrava buono. Mi girai alla ricerca di un tavolo e quasi mi schiantai
col vassoio contro quello di Shauna.

«Ehi! Ehm, Shauna, giusto?», chiesi nel modo più noncurante possibile.

«Ehi, tu.» Guardò il mio vassoio. Sul suo c'erano solo un sandwich, i
fagiolini e una mela. «Vuoi prendere un tavolo?»

Lo volevo? Non esseno certo di riuscire a parlare senza dire qualcosa di
stupido, mi limitai ad annuire e la seguii a un tavolo che non era
troppo affollato. «Sai cosa stavo facendo il Giorno di Colombo?», dissi
in modo da aprire una conversazione, solo che avrebbe potuto suonare
come un modo per vantarsi, così aggiunsi: «Hai fatto qualcosa di stupido
quel giorno?».

«Ero con amici a Phoenix.»

«Non eri in giro a guardare le foglie quel giorno?»

Inclinò la testa e mi lanciò un'occhiata sarcastica. «La gente a Phoenix
non sta tanto a guardare le foglie.»

«Capito. La mia famiglia non ha bisogno di andare a cercare le foglie,
nel Giorno di Colombo ce n'è uno strato di mezzo metro sul prato.»

«Non mangi?» Guardò il mio vassoio ancora pieno.

«Oh, sì. Non voglio parlare a bocca piena.» Diedi un morso a un sandwich
e mandai giù in fretta.

«Sono andata prima a casa per assicurarmi che mia madre stesse bene, mio
padre era in viaggio d'affari a Houston, poi sono andata a controllare
mia nonna. Vive in un grattacielo e, senza elettricità, era bloccata
lassù, non pensavo che sarebbe riuscita a scendere le scale, ha un'anca
malconcia. Ed è, sai, vecchia.» La voce le tremava un po', ricordare
quel giorno stava suscitando in lei molte emozioni. «Era così
spaventata, era così spaventata e mi disse di non preoccuparmi per lei,
mi disse che ero giovane e forte e che dovevo uscire dalla città, andare
in un posto dove gli alieni non mi avrebbero trovata. Continuava a dire
che era la fine del mondo, il giorno del Giudizio universale e che
dovevo lasciarla e sopravvivere. In città, non c'erano navi d'assalto
ruhar che atterravano per la strada, tutto ciò che vedevamo erano luci e
scie nel cielo e colonne di fumo provenienti da siti che avevano
colpito. L'unica cosa che ci informava del fatto che era in corso
un'invasione aliena erano gli annunci di emergenza alla radio. E
all'inizio non ci credevo.» Mangiò l'ultima forchettata di fagiolini,
poi spinse via il vassoio. «Quei criceti hanno spaventato mia nonna
quasi a morte. Li odio per questo. Ecco perché sono qui fuori. Non
voglio che mia nonna, o chiunque altro, debba mai più avere così tanta
paura.»

«Sì. Prima uscivo nelle notti limpide, guardavo il cielo e mi chiedevo
cosa ci fosse là fuori nell'universo. Dopo il Giorno di Colombo guardo
il cielo e, se le stelle scintillano per un secondo, penso che possa
essere una nave nemica che salta in orbita.»

«Sì! Esatto!» Colpì il tavolo col palmo della mano. «Per colpa dei
criceti, nessuno può più guardare il cielo senza paura. Odio quei figli
di puttana per questo.»

«Ci hanno rubato l'innocenza?», mormorai con la bocca piena di purè di
patate.

«Sì, qualcosa di simile. Non a te e a me, noi abbiamo visto
combattimenti, di merda ne abbiamo vista in Nigeria. Ma la maggior parte
della gente, che vive la sua vita e basta, non riuscirà a pensare di
essere di nuovo al sicuro, mai più.» Distolse lo sguardo per un istante.
«Mai più. Adesso sappiamo che c'è un'intera galassia piena di nemici qua
fuori. Che incombono sulle nostre teste, sempre.»

Parlammo di che grandissima stronzata fosse aspettarsi che andassimo in
guerra con la stessa attrezzatura che l'esercito degli Stati Uniti aveva
usato per decenni e condividemmo voci di corridoio rispetto a quello che
sarebbe stato il nostro primo dislocamento, nessuno di noi sapeva nulla
rispetto a dove i kristang ci avrebbero mandato. E ci chiedemmo
preoccupati che cosa stesse succedendo a casa, mentre io finivo la cena
e lei mangiava la mela. La sua famiglia se la passava molto peggio
rispetto ai miei genitori: i suoi avevano entrambi perso il lavoro e,
con il suo giovane fratello e la nonna, si erano trasferiti nel Texas
orientale, dove il padre era in una squadra di lavoro agricolo e la
madre lavorava in un asilo nido. Era stato un grande cambiamento per i
suoi genitori, la madre era stata agente immobiliare e il padre
rappresentante delle vendite per un'azienda informatica. Entrambi i
posti di lavoro erano svaniti quando l'economia era crollata e il loro
appartamento di Phoenix era diventato invivibile senza elettricità per
far funzionare i condizionatori.

Mio padre aveva perso il lavoro alla cartiera, ma per i primi due mesi
se l'era cavata con lavori saltuari, taglio e trasporto di legname,
saldatura e riparazione auto. Mia madre aveva ancora un posto fisso come
insegnante, specie dopo che così tante famiglie si erano trasferite in
campagna da città ormai prive di elettricità. Alla fine di dicembre, a
qualcuno era venuto in mente che la grande cartiera dove mio padre aveva
lavorato, sede di un impianto di cogenerazione che in pratica bruciava
segatura avanzata per produrre elettricità per il mulino, avrebbe potuto
essere reimpiegata nella produzione di elettricità per la regione.
Questo aveva fatto tornare mio padre a lavorare al mulino, con l'unico
problema che non avevano abbastanza legname da portarci, perché non
c'era gasolio a sufficienza per i camion destinati al trasporto della
legna. Nei Great North Woods c'era un sacco di legname, ma non c'era
modo di portarlo da nessuna parte. In giro si trovava a stento
carburante diesel a sufficienza per i treni. La soluzione era stata
portare vecchie locomotive a vapore da una ferrovia turistica nel New
Hampshire, con ragazzi che ancora sapevano come farle funzionare e come
alimentarle a legna, piuttosto che a carbone o petrolio. Buon vecchio
ingegno americano in azione, con treni a vapore che trasportavano i
tronchi dal bosco a una cartiera, diventata una centrale elettrica.
L'intera rete elettrica degli Stati Uniti era stata ricostruita,
lentamente, come la ``rete intelligente'' che avremmo dovuto avere
decenni prima, con un sacco di piccole centrali elettriche locali,
piuttosto che un impianto gigante per un'intera regione. I kristang non
erano stati di grande aiuto nel ripristino della produzione di energia
elettrica in tutto il pianeta, tranne che nelle aree industriali
ritenute vitali per lo sforzo bellico. I nostri nuovi alleati
disprezzavano il fatto che l'umanità si affidasse ancora ai combustibili
fossili per l'energia, e avevano installato giusto un paio di reattori a
fusione sulla Terra per i propri bisogni, ma, anche se avessero spiegato
la loro tecnologia di fusione ai nostri scienziati, ci sarebbe voluto
molto tempo prima che l'umanità fosse in grado di costruire da sola un
reattore a fusione funzionante.

Quando ero partito, l'elettricità era ancora scarsa nella mia città
natale, la maggior parte dell'energia della nuova centrale elettrica di
Milliconack andava a valle della zona di Bangor. I miei genitori
riscaldavano la casa con stufe a legna, non avevano bisogno di aria
condizionata e il giardino, che avevamo ampliato durante l'inverno,
avrebbe prodotto un sacco di raccolto per sfamare la mia famiglia e le
persone che vivevano con lei, oltre a cibo in eccesso destinato alla
vendita. Avevano due mucche e galline per le uova e, quando ero partito,
mia madre era in trattativa per un paio di maialini. Era così in tutto
il mondo: le persone che erano più in alto sulla scala economica prima
dell'attacco dei ruhar ora, in certi casi, stavano peggio di quelle che
già prima tiravano a campare. Se eri un contadino prima del Giorno di
Colombo, la tua vita non era cambiata molto -- a parte il fatto che era
difficile trovare carburante per i trattori -- dato che la gente aveva
ancora bisogno di mangiare, e quindi c'era un sacco di domanda. Cosa
dice la Bibbia, qualcosa come: ``Molti dei primi saranno ultimi e molti
degli ultimi saranno i primi''? Nessuno si aspettava che le profezie
delle Scritture si sarebbero realizzate per via di un'invasione aliena.

«Mi sto facendo il culo per qualificarmi al servizio attivo», spiegò
Shauna. «Quasi qualificata quando mi sono iscritta, sono caduta su un
percorso a ostacoli e mi sono slogata il legamento crociato anteriore a
due settimane dalla fine. Mi ha fatto incazzare! Così sono passata alla
logistica. Volevo il servizio attivo, Joe.»

Essendo in fanteria ed essendo stato in servizio attivo, volevo metterla
in guardia contro l'eccesso di entusiasmo. Come voleva servire non era
affar mio, così tenni la bocca chiusa. Finimmo di mangiare, la gente
della sua unità venne a parlare con lei e poi se ne andò. Se questo
poteva essere considerato un primo appuntamento, fu deprimente.

\hfill\break

La prima simulazione di guerra su Campo Alfa fu l'operazione Scudo
valoroso. I nomi in codice per le operazioni dovevano idealmente essere
scelti a caso, in modo da non dare al nemico un'idea del tuo intento.
``Operazione Tempesta nel deserto'', per esempio, era un nome figo, ma
gli iracheni non avevano avuto difficoltà a capire che si trattava di
un'offensiva nel deserto. Non che facesse differenza allora. Il problema
con le parole casuali è che, senza volerlo, si potrebbe finire per avere
operazioni dai nomi divertenti, o claudicanti, tipo ``Coniglio
fiammeggiante'' oppure ``Giustizia zoppa'', perciò in realtà i nomi
delle operazioni erano inventati, o almeno controllati, dallo staff
delle pubbliche relazioni dell'Unef. È probabile che il nome ``Scudo
valoroso'' fosse stato pensato per ispirare le truppe, ma quando lo
sentii, l'unica cosa che mi venne in mente era che dovevamo morire in
modo valoroso. Come in quel vecchio detto romano, per cui un soldato
doveva tornare a casa ``o con lo scudo, o sopra lo scudo''.

Vorrei potervi dire che durante la nostra prima simulazione di guerra
inventai una tattica incredibilmente brillante, a cui nessun altro aveva
pensato, e che ci portò, o almeno portò il mio plotone, alla vittoria.
Potrei dirvelo e sarebbe una storia fantastica, giusto? Forse un giorno
racconterò qualcosa del genere ai miei nipoti, quando si lamenteranno di
quanto il loro vecchio nonno sia fuori moda. O forse racconterò una
storia del genere un giorno, quando avrò alzato il gomito e penserò che
mi possa far rimediare una scopata. La verità, purtroppo, è che io e
tutti gli altri nei due plotoni della nostra compagnia fummo dichiarati
morti, a causa di un attacco orbitale, meno di un'ora dopo l'inizio
della simulazione di guerra. Stavamo sfrecciando in un canyon poco
profondo, passando di copertura in copertura, quando il mio zPhone suonò
e una voce stridula annunciò che ero morto. Poi il dispositivo smise del
tutto di funzionare. Merda. Come cazzo si fa a sprecare un colpo di
cannone a rotaia per due plotoni di fanti? La compagnia si era divisa
per evitare di fornire un bersaglio succoso, ma non aveva funzionato.

Il protocollo ci chiedeva di sederci o sdraiarci nel punto in cui ci
trovavamo e di aspettare un segnale chiaro, che sarebbe potuto arrivare
il giorno dopo a seconda di come sarebbe andata la simulazione. Il sole
era già sopra l'orizzonte e faceva già caldo, i nostri ufficiali
decisero che potevamo imbrogliare un po' e tornare indietro attraverso
il canyon per ripararci sotto uno strapiombo di roccia. È lì che passai
il resto di quella giornata, tutta la notte e metà della mattina
successiva. Eravamo a corto d'acqua, eravamo a sei ore di cammino dalla
base, perché ci eravamo lanciati dalle Poiane prima dell'inizio della
simulazione di guerra. Qualcuno al quartier generale ebbe pietà di noi e
ci mandò dei camion, non avevano l'aria condizionata ma, anche se
brontolammo per questo, apprezzammo il fatto di non dover scarpinare per
tutto il percorso.

\hfill\break

Una lezione che la simulazione ci stampò a fuoco nella mente è che la
chiave per sopravvivere al combattimento a terra in quel tipo di guerra,
in una situazione in cui la posizione strategica è sopra l'atmosfera, è
la stessa di quando devi sopravvivere al combattimento in situazioni in
cui il nemico potrebbe usare armi nucleari tattiche. Evitare di
concentrare le forze e mantenere la copertura il più possibile. In una
potenziale situazione di guerra nucleare, non date al nemico un
bersaglio che potrebbe essere tentato di colpire con una testata
nucleare, come forze dalle dimensioni di un battaglione, grandi depositi
di munizioni, campi d'aviazione fissi, ponti vitali, cose del genere. I
kristang non usavano armi nucleari contro obiettivi a terra: si potrebbe
pensare che fosse una buona cosa, ma in realtà rese la situazione più
pericolosa. Decidere di usare una testata nucleare comporta un sacco di
ostentata preoccupazione della leadership di alto livello per
l'eventualità di un'enorme escalation nucleare del conflitto,
un'escalation da cui non si potrebbe fare un passo indietro, un genio
che non si potrebbe rimettere nella lampada. L'uso di cannoni a rotaia,
che possono accelerare i colpi fino a raggiungere una percentuale
significativa della velocità della luce e che possono liberare
l'equivalente di energia distruttiva di una testata nucleare tattica, è
così facile da essere scontato nella guerra interstellare. Per noi
soldati semplici a terra, questo significa che essere colpiti da un
cannone a rotaia dall'orbita era più una questione di quando, che di se.

Nel complesso, nella simulazione, le forze Blu dell'Unef furono
schiacciate dalla forza Rossa simulata dei ruhar. Si dice che i kristang
fossero molto soddisfatti del risultato. Non soddisfatti che i soldati
delle truppe dell'Unef fossero ``morti'' a migliaia, ma soddisfatti che
le truppe dell'Unef avessero resistito abbastanza a lungo da ritardare
l'occupazione della superficie da parte delle forze Rosse. Le forze
sparpagliate Unef avevano messo in atto tattiche di guerriglia per
attaccare le forze Rosse dopo il loro atterraggio e le nostre squadre
armate di Zinger Manpad avevano ``abbattuto'' un numero incoraggiante di
navicelle Rosse. E anche i nostri aviatori si erano comportati bene.
Certo, quando i kristang fermarono la simulazione di guerra, la forza
Blu dell'Unef aveva meno di cento aerei in grado di volare. Se pensate
che sia un male, sappiate che i kristang si aspettavano che tutte le
nostre risorse aeree fossero esaurite entro le prime dodici ore. Il
combattimento moderno è ad alta intensità, specie una volta che i piedi
si sono staccati da terra.

Due giorni dopo la simulazione di guerra, la sera, i componenti del mio
gruppo di fuoco stavano guardando una partita di basket, il nostro
battaglione contro una squadra dell'esercito britannico, ci stavamo
godendo hot dog e popcorn, quando un tenente dell'esercito ancora in
divisa di volo si sedette accanto a noi in tribuna. «Ehi, lei è un
pilota, signore?», chiesi, stupidamente.

«Sì, sì», disse l'altro mordendo un hot dog, «un tempo pilotavo un
Apache, ora un Pollo.»

«Com'è andata nella simulazione?», domandai con entusiasmo. Avevo
sentito raccontare molto dai fanti, ma niente dai nostri aviatori fino a
quel momento. Nel mio entusiasmo, non presi in considerazione il fatto
che il tipo potesse avere appena finito un lungo volo e volesse
rilassarsi e guardare una partita piuttosto che raccontarmi la stessa
storia che aveva già raccontato un centinaio di volte. Non c'era motivo
di cui preoccuparsi, perché un pilota non si stanca mai di parlare del
volo.

«Me la sono cavata bene», disse guardando la partita, «ho abbattuto un
Dodo e un Avvoltoio.»

«Cazzo, signore, lei è un grande!» Gli offrii di darmi il cinque e fui
entusiasta quando lo batté.

«Grazie, soldato. Ovvio che sono stato colpito dal missile di un
incrociatore ruhar due minuti dopo. Mi sentivo bene comunque, ho vissuto
più a lungo di quanto mi aspettassi.»

«Ha abbattuto un Avvoltoio?», chiese Ski. «È un velivolo d'assalto?»

«Sì. I Dodo avevano fatto sbarcare le truppe, una formazione di sei Dodo
con Avvoltoio di scorta, stavano tornando in orbita quando abbiamo
attaccato da due direzioni, è arrivato rapido e radente, proprio sul
ponte. È stato un combattimento aereo imponente, uno scontro durissimo,
abbiamo preso tutti i Dodo, simulati, ovvio, e abbiamo perso metà dei
nostri aerei. Quell'ultimo Avvoltoio, poi\ldots{} Ho dovuto inseguire
quella cosa puzzolente sopra i ventunomila metri, che è l'altitudine
massima per un Pollo, non siamo autorizzati a operare fuori
dall'atmosfera, ero in massa di reazione interna ed ero a stento in
grado di controllare la nave con i propulsori secondari. L'Avvoltoio era
molto sopra di me, stava spingendo a 2 g e saliva allontanandosi, mentre
io stavo per perdere il controllo e andare nel panico. Ho lanciato a
raffica i miei ultimi quattro missili e il mio mitragliere ha colpito
l'Avvoltoio con il raggio di particelle, il raggio è solo difensivo, ma
ha fiaccato le difese dell'Avvoltoio abbastanza da far passare un
missile.» Si voltò per guardare la partita. «Il loro stealth attivo non
funziona bene fuori da un'atmosfera, non ce l'hanno detto, ma l'ho
capito.» Prese un'altra manciata di popcorn dal sacchetto di carta e li
mangiò lentamente, pensoso. «Sai, il cielo a quell'altitudine non è blu
brillante, è quasi nero, e puoi vedere che il pianeta che si estende
sotto di te è rotondo. Quell'incrociatore ruhar che mi ha preso? Potevo
vederlo. Non come un fantasma sul visore a sovrimpressione per la
simulazione, questa era una vera nave, proprio lì, sospesa nel cielo.
Era inquietante. L'immagine sul visore mostrava una nave ruhar
sovrapposta alla vera nave kristang per la simulazione, quando ho spento
il visore, ho potuto vedere la vera nave. Una vera nave stellare, in
orbita bassa sopra di me. È stato fantastico, per un momento, finché non
ha colpito la mia gondola di dritta con la punta di un laser. Ci sono
volato dritto dentro. Mi ha fatto vorticare, ero occupato a cercare di
tenere insieme la nave, quando il loro missile mi è strisciato su per il
culo. Simulato, sembrava abbastanza reale. Non sarei neppure stato in
grado di recuperare, il pilota automatico ha preso il sopravvento e mi
ha guidato fino ad angeli trenta, cioè trentamila piedi.»

«Merda», disse Ski, «pensavo che noi fanti a terra ce la fossimo vista
brutta.»

«È così, senz'altro. Noi siamo più esposti in cielo, abbiamo
contromisure che aiutano molto con le minacce missilistiche. È buffo, è
un po' come fossimo tornati ai tempi della seconda guerra mondiale,
prima che i missili guidati la facessero da padroni nel combattimento
aereo. Lo stealth attivo e le contromisure significano ottenere una
soluzione di fuoco affidabile sul nemico, non è una chiusa come nel
combattimento aereo sulla Terra: un missile colpisce un bersaglio forse
il 15, 20\% delle volte, i kristang non ce lo danno per certo.
Immaginiamo il 20\% in base alle tattiche che ci hanno insegnato.» Rise.
«Ridicolo: mi hanno mandato qui perché pensavano che per un pilota di
elicotteri la transizione per pilotare un Pollo o una Poiana sarebbe
stata più facile. Quando sono in volo a punto fisso, certo, è abbastanza
vero. Il fatto è: un Pollo può andare a velocità supersonica e salire
più in alto dei nostri jet. Non ero pronto per questo; l'ultima volta
che ho volato ad ala fissa era su una piccola turboelica in
addestramento. Il mio Pollo può salire e superare facilmente un F22, per
la maggior parte del tempo che lo stai guidando, la cosa è un jet, non
un elicottero. Una Poiana è più simile a un V22 Osprey, non è
supersonica.» Scosse la testa. «Le tecniche che abbiamo imparato
guidando aeromobili ad ala rotante sulla Terra non si applicano qui,
devo continuare a ricordarmi di non preoccuparmi di stress dinamico
sulle pale del rotore e non ho bisogno di tenere il rotore di coda
libero quando atterro, perché non ce l'abbiamo.»

«Signore», chiese Cornpone, «noi abbiamo solo aerei, già, e il nemico ha
navicelle che possono andare nello spazio. Abbiamo davvero una chance
nel combattimento lassù?» Indicò il cielo.

«Oh, sì. Ho capito cosa intendi. Sì, ce l'abbiamo, abbiamo davvero un
vantaggio sulle navicelle nel combattimento aereo. Le navicelle sono
grandi e pesanti, si muovono come bisonti quando sono in profondità in
un'atmosfera. Con stealth e contromisure, se la tua prima raffica di
missili non va a segno, chiudi la distanza dal bersaglio molto in fretta
e poi entri in fase di manovra tattica e utilizzi le tue armi. La chiave
per il combattimento aereo è l'energia cinetica, se perdi troppa
velocità in una volta sei morto. Energia significa che puoi ingaggiare e
disinnescare uno scontro aereo secondo la tua necessità. Le navicelle
sono alla fine molto più veloci dei nostri aerei, possono raggiungere la
velocità di fuga, ma non accelerano altrettanto in fretta. Noi possiamo
accelerare molto più in fretta.» Continuò a parlare ancora per un po'
delle tattiche di combattimento aereo, noi tre eravamo entusiasti, in
quel momento pensai che fosse il tipo più figo mai esistito. Lasciai il
mio posto e gli andai a prendere un altro hot dog, così avrebbe
continuato a parlare. La partita finì, però; la squadra britannica vinse
per due punti e arrivò un gruppo di piloti, così lui se ne andò.

«Accidenti», disse Ski con fare malinconico, «vorrei essere un pilota.
Ci ho pensato anch'io, uno dei miei zii è in un aeroclub a casa, il club
ha un solo motore, ehm, Piper, qualcosa del genere. Una volta, a
Minneapolis, mio zio portò in volo me e mio padre per una partita dei
Bears, fu grandioso. Mai avuti i soldi per le lezioni di volo, però.»

«Sarebbe fantastico», concordai.

«Sì, abbastanza. Comunque è morto nella simulazione», sottolineò
Cornpone.

«Sì, sì, oh, sì», disse Ski. «Anche noi, e non abbiamo fatto un cazzo
prima di morire. Preferirei morire in cielo facendo qualcosa, piuttosto
che essere vaporizzato da un colpo di cannone a rotaia che non mi
aspettavo.»

\hfill\break

Ispirati dalla visione della partita di basket, Ski, Cornpone e io
andammo in un altro campo e ci imbattemmo in una partitella improvvisata
cui potevamo unirci. Ski aveva un ottimo tiro esterno, mi concentrai sul
passargli la palla e andare a rimbalzo in difesa. Era bello giocare
duro, sfogare la frustrazione col sudore. Tutti noi eravamo incazzati
perché eravamo morti e non avevamo ottenuto nulla durante l'operazione
Scudo valoroso. Certo, farci saltare in aria dall'orbita era uno
scenario realistico, la nostra domanda era: quale insegnamento avremmo
dovuto trarne? Inoltre, come avrebbero fatto i kristang a valutare le
nostre capacità di combattimento, se non avevamo fatto altro che
ottenere un lancio di dadi sfortunato da un computer kristang in orbita?
Se quei dadi immaginari fossero rotolati in modo diverso, qualche altro
sfortunato bastardo sarebbe morto al posto nostro.

Lo zPhone alla mia cintura vibrò durante la partita, non lo controllai
finché non facemmo una pausa per bere. Era un messaggio di Shauna,
voleva sapere se fossi impegnato e, in caso contrario, se volessi uscire
con lei.

Volevo?

Il sole sorgeva a est? Sulla Terra, intendo.

Certo che sì!

Ecco un semplice test per verificare se un ragazzo vuole incontrarsi con
una ragazza.

Primo step: ha polso?

Secondo step: è cosciente?

Non c'è bisogno di un terzo step.

Da idiota che sono, digitai: «È una proposta indecente?». Emoticon
sorridente con occhiolino.

Per fortuna, il mio istinto di conservazione m'indusse a cancellare il
messaggio prima di premere ``Invio''. Invece, digitai, col massimo della
noncuranza che mi riusciva: «Certo! Sto giocando a basket, stiamo
vincendo. Finisco presto».

E poi mi presi degli accidenti da Cornpone e Ski perché era chiaro che
non c'ero più con la testa, continuavo a immaginarmi vibrazioni di
notifica di messaggi in entrata sul mio zPhone. Vincemmo la partita
comunque, non che qualcuno facesse sul serio tenendo il punteggio.

Cazzo. La partita era finita e non avevo ricevuto risposta da Shauna.
Ero stato troppo impaziente? Troppo noncurante? Merda. Mia nonna
raccontava storie di come se ne stesse seduta accanto al telefono a casa
sua, questo avveniva prima dei cellulari, seduta accanto al telefono,
dicevo, in attesa che un ragazzo chiamasse. Mi sembrava patetico,
allora. Seduta accanto al telefono, invece di essere fuori a divertirsi,
nella speranza che un perdente decidesse di chiamare?

Eppure, quando Ski e Cornpone suggerirono di andare a vedere che film
stavano proiettando nella grande tenda del centro ricreativo, invece di
andare con loro, dissi che ero stanco e volevo dormire. Mi guardarono in
modo strano, ma mi lasciarono in pace. Ed eccomi lì, non con i miei
amici, non di ritorno alla nostra tenda ad aspettarli, nel caso ci
fossero tornati, ero lì, seduto da solo nella tenda del campo da basket
ad aspettare che una ragazza mi contattasse.

Per tutte le donne che hanno aspettato che un ragazzo chiamasse o
mandasse un messaggio: so come vi sentite. Fa schifo.

E poi il mio zPhone notificò la ricezione un messaggio.

\hfill\break

Quello che mi aspettavo era che Shauna suggerisse di vederci da qualche
parte, parlare, magari incontrare i suoi amici, cose così. Mi spiazzò
del tutto chiedendomi d'incontrarci a un cancello del parco macchine del
gruppo logistico, un'area recintata per i camion kristang che stavamo
usando. Aveva il codice del cancello ed entrammo, la seguii in silenzio
mentre si metteva un dito sulle labbra. Facemmo lo slalom tra i camion
parcheggiati, non sapevo dove mi stesse portando, finché non ci fermammo
dietro un camion e lei chiuse il cancello sul retro. Il camion aveva una
specie di tetto di tela, lei tirò il lembo da una parte e puntò una
torcia all'interno. Per un secondo, ebbi paura che ci fossero altre
persone nel camion e che Shauna si aspettasse che prendessi parte a un
qualche losco traffico che non mi interessava: il mio fascicolo
personale dell'esercito conteneva già abbastanza osservazioni su tutte
le volte che mi ero messo nei guai.

Non c'era nessuno sul camion. Solo uno spesso mucchio di coperte.

I miei occhi dovevano aver mostrato la mia completa sorpresa, perché
Shauna mi afferrò per la maglia e mi baciò con passione. «Ti va bene?» I
suoi occhi brillavano tra le fioche luci di sicurezza. «Sei carino. Un
po' stupido, forse, ma carino.»

Feci di sì con la testa come un bambino di quattro anni quando gli
chiedono se vuole una ciotola di caramelle. «Sì, ok. Sì!»

\hfill\break

Per quelli tra voi che hanno fatto sesso, non vi annoierò con dettagli
sudaticci su di me e Shauna. Sapete com'è, e sono sicuro che non abbiamo
fatto niente che voi non abbiate fatto prima. O abbiate desiderato fare.

Per quelli tra voi che non l'hanno mai fatto, non vi voglio rovinare la
sorpresa.

Indizio: f-a-n-t-a-s-t-i-c-o.

Wow.

\hfill\break

Ci prendemmo una pausa, una chiacchierata, Shauna stava sdraiata sulla
schiena, io accanto a lei. Dopo un po', le feci scivolare con dolcezza
le dita giù per il collo, tra i seni e lo stomaco, e giù ancora.

«Che cosa stai facendo?», ridacchiò.

«Ti sto facendo passare le dita nella chioma», dissi con fare innocente.

«Non è così che stanno le cose!» Rise e mi schiaffeggiò per scherzo.

«Non è romantico?»

«Sì, come il tuo amico qui», e tamburellò con le dita il mio amico. «È
tutta una questione di romanticismo, Joe.»

«Ehi», dissi mentre il mio amico si svegliava al tocco delle sue dita,
«lui è molto romantico.»

«Mmh, sì, ci scommetto. È romanticissimo, se serve a fare sesso.»

«Ok, devo ammetterlo, è un cane in calore, non riesco a controllarlo a
volte.» Proprio come in quel momento, con Shauna che continuava a
incoraggiarlo.

«Oh, non hai il controllo su di lui? Come se tu non c'entrassi niente?»

«Non ne hai idea», scossi la testa con mestizia. «A volte mi sveglio
alle quattro del mattino perché lui ha dimenticato la chiave della porta
d'ingresso e io gli faccio: ``Dove sei stato?'', e lui mi fa: ``Da
nessuna parte, sono solo uscito per una passeggiata''. E odora di vodka
e profumo, e io so solo che è andato in giro a cacciarsi nei guai.»

«Una chiave? E dove la metterebbe? Esce da solo?» Rise.

«Beh, si porta dietro i ragazzi, quei tre sono una banda. Ma, è chiaro,
io non ho colpa dei casini in cui si cacciano. Il mio cuore è puro.»

«Il tuo cuore è una pura stronzata, Joe, ma sei divertente.» Rise di
nuovo. «Parlerò con il tuo amico, vediamo cosa ne pensa. Non ti
dispiace, vero?»

«No, no, per favore, continua, tu puoi parlare con lui quanto vuoi.»

\hfill\break

Un paio d'ore dopo, Shauna disse che dovevamo andarcene, perché stava
arrivando una pattuglia di sicurezza. Di solito accendeva le luci solo
per un minuto, per sincerarsi che non ci fosse nulla di grave. Capivo il
loro punto di vista, chi avrebbe rubato un camion su quel pianeta? Che
cosa se ne sarebbe fatto?

«Ehm, allora ti chiamo? Domani?» Cazzo. Avrei dovuto dirlo prima di
rimettermi la maglia e i pantaloni. In mia difesa, nel camion stava
cominciando a fare freddo, ora che ero fuori dalle coperte.

«Chiamarmi? Tu mi chiamerai?»

«Oppure ti scrivo?», dissi lentamente, col sospetto di avere fatto
qualcosa di sbagliato.

«Tu chiamerai me? Ehm.» Era contrariata da qualcosa. «Vorrei essere
incazzata con te, ma sei così maledettamente stupido e carino.»

«È vero, è una maledizione. Sai, non voglio che tu pensi che io sia,
ehm, sai?»

«Cosa? Che mi stai usando per il sesso? Joe, non sono una stupida
liceale. Sono io che sto usando te per il sesso.»

«Ah.»

«È un problema?»

«No! No, non è affatto un problema», mentii. Era, forse, un piccolissimo
problema, per il mio ego. Sì, lo so, questa sarebbe una situazione da
sogno per la maggior parte dei ragazzi, a meno che non siano stati del
tutto onesti con se stessi. A un ragazzo piace sapere che ha significato
qualcosa per una ragazza, anche se lei è stata una botta e via.

Sì, è una questione di ego maschile. Fatemi causa.

«Senti, Joe, siamo su un pianeta alieno», disse lei con la voce
soffocata dalla coperta sopra la testa, contorcendosi per entrare nei
pantaloni mentre era ancora bella e calda sotto la coltre. Vorrei averci
pensato. «Non sappiamo dove finiremo, potremmo essere inviati su pianeti
diversi la prossima cazzo di settimana. Mi sto facendo il culo per
qualificarmi al servizio attivo, voglio fare la differenza in questa
guerra, questo è quello su cui mi sto concentrando ora. Io non ho tempo
per un fidanzato. Tu non hai tempo per una fidanzata. Mi piaci, ti
piaccio, le ragazze si eccitano e sei abbastanza bravo a letto. Nessun
impegno. Possiamo evitare complicazioni?»

«Oh, ehm», la mia giornata era tornata a essere buona, adesso che
capivo, «certo. Sì, sì!»

«Grandioso.» Lei fece dondolare le gambe fuori dalla coperta, infilò i
piedi negli stivali, saltò in piedi e mi baciò sulla guancia. «Sarò io a
chiamare te.»

Questo mi spaventò.

Poi, chiamò due sere dopo e consumammo le sospensioni del camion.

Fu fantastico.

\hfill\break

Ogni illusione che le mie avventure amorose notturne fossero passate
inosservate finì fuori dalla finestra a colazione la mattina dopo. Il
Sergente Koch ci aveva fatto correre fino allo sfinimento su un percorso
a ostacoli nel buio, io non avevo dormito abbastanza e tutti e tre
avevamo una fame da lupi.

«Cos'è quello?», mormorò Cornpone a bocca piena, indicando la scodella
di Ski con un pezzo di pane tostato.

Lo guardai anch'io. «Cazzo, Ski», dissi, «ma che fai?» Invece di uova,
french toast, frittelle di patate o qualsiasi altra cosa buona, aveva un
solo pezzo di pane tostato appena imburrato e una scodella di quello che
sembrava essere porridge semplice. Semplice come l'avena inzuppata
nell'acqua. Senza uvetta, senza zucchero di canna, senza noci, nemmeno
una crema. Il genere di cose che servono ai cavalli. Già. «Sembra una
scodella di, tipo, zuppa\ldots{} tristezza.»

«Una scodella di tristezza.» Cornpone rise fino a soffocare, sputando
uova nel piatto. «È divertente. Zuppa di tristezza è un marchio della
Kellogg's?»

«Sul serio», aggiunsi, mentre Ski se ne stava seduto là con la faccia
impietrita, «mangiare quella roba è come essere un fan dei Chicago Cubs.
Certo, una volta ogni tanto i Cubs avranno la fortuna di vincere, ed è
come trovare una toffoletta nella tazza, ma andiamo, tu lo sai che prima
della fine della stagione soffocherai in una scodella di tristezza.»

«Taci.» Ringhiò Ski, ma gli brillavano gli occhi.

Cornpone colpì il tavolo con la mano. «Bish, cazzo, mi stai uccidendo.»

Offrii una fetta di french toast a Ski, ma lui la sventolò con
un'espressione accigliata. «No, amico, quel polpettone l'altra notte ha
litigato col mio stomaco, ho bisogno di qualcosa di leggero dopo la
corsa a ostacoli.»

«Ah, sì, il polpettone. Cos'era quella sbobba?» Cornpone scosse la
testa. «Olio di motore 10W-40, almeno dal sapore.»

«No, penso che fosse il fluido idraulico che i criceti usano per quelle
Poiane. Squisito», mi leccai le labbra, «carne misteriosa che galleggia
in una pozza oleosa di fluido idraulico.»

«Sì, sì», Cornpone concordò, cacciandosi in bocca patate fritte a
cucchiaiate. «E quando la salsa oleosa inizia a raffreddarsi, si forma
quella pellicola gommosa in cima, e lei\ldots»

«Ah.» La faccia di Ski stava diventando verde. Spinse via il porridge.
«Ragazzi, zitti, per favore, non ce la faccio più.»

Rompersi le palle a vicenda era elemento essenziale dell'essere un
gruppo di fuoco, sapevamo anche quando smettere. L'ultima cosa che io e
Cornpone volevamo era che Ski vomitasse su tutto il tavolo della
colazione, cosa che ci avrebbe reso molto popolari nel plotone. No.
«Ecco», presi due fette di pane integrale da Cornpone e le misi sul
piatto di Ski, «questo ti riempirà. Abbiamo una mattinata impegnativa.»
Il programma prevedeva un'ora al poligono di tiro, poi il plotone
avrebbe fatto un'escursione di venti chilometri, raggiungendo cinque
punti di navigazione prima della pausa pranzo. Questo era parte della
nostra preparazione per una grande simulazione di guerra di armi
combinate che ci sarebbe stata di lì a una settimana, quando speravamo
di mostrare ai kristang che soldati cazzuti eravamo noi umani nel
combattimento simulato. Parte dello scenario della simulazione di
guerra, che era ancora in fase di sviluppo, includeva il bombardamento
orbitale simulato e prevedeva esseri umani che combattevano contro
navicelle kristang alla guida di aerei sottratti ai ruhar. L'Unef
sarebbe stata la squadra Blu che difendeva il pianeta e i kristang
sarebbero stati la squadra Rossa degli invasori, con la forza kristang
che faceva la parte dei ruhar, usando tattiche dei ruhar.

Un paio di fette di pane nello stomaco di Ski lo fecero sentire meglio,
al punto da fargli rubare un pezzo di french toast dal mio piatto. Si
sentiva bene abbastanza da chiedere: «Ehi, Bish, quella ragazza, Shauna,
ci stai dando dentro?».

«Sì», disse Cornpone e mi diede una gomitata nelle costole, «come stanno
le cose?»

Per colpa loro quasi soffocai con un boccone di patate fritte. «Le cose?
Le cose stanno così, anzitutto è una donna, non una ragazza.»

«Sai cosa intendo.» Cornpone non aveva intenzione di sviare
dall'argomento.

«E non sono affari che vi riguardano. Ci siamo scontrati a cena e
abbiamo parlato. Non è una gran cosa.»

«Hai parlato con una ragazza vera, viva, reale. Potresti avere toccato
una ragazza vera. Questa è una gran cosa», obiettò Ski.

E Cornpone aggiunse: «Bish, lo sai che c'è? Ho sentito che in tutta
l'Unef solo il 16\% delle forze è costituito da donne, il 16\%! Questo
significa che per ogni ragazzo, ci sono tipo, ehm, per ogni ragazza ci
sono tipo\ldots{} quattro ragazzi!».

«Penso che tu debba controllare i tuoi calcoli, Jesse.»

«Il punto, Bish, è che con queste probabilità siamo un po' disperati, e
se un tizio a caso su questo pianeta sta scopando, vogliamo almeno dei
dettagli.»

«Non ho intenzione di darvi dettagli.»

«Ah ah!» Cornpone diede una pacca sul tavolo, attirandosi gli sguardi di
altri attorno a noi. «Ci sono dettagli.»

Merda. Avevo rovinato tutto. «Ragazzi, mettiamo il caso che un tizio su
questo pianeta si stia, diciamo, godendo la compagnia di una signora.
Pensate che le probabilità di altri ragazzi di avere fortuna
aumenterebbero se ne parlasse o se tenesse la bocca chiusa? Eh? Pensate
che le donne su questo pianeta sarebbero felici se un tipo si vantasse
dei dettagli?»

«Cazzo», brontolò Ski, «ha ragione.»

«O ti vanti di scopare, o scopi per davvero. È molto semplice»,
conclusi.

«Cazzo. Diavolo, in questo caso», disse Cornpone scontroso e tese la
mano con la forchetta, «mi prendo il resto del tuo french toast.»

\hfill\break

Una mattina, il mio zPhone suonò mentre ero al poligono di tiro, era un
messaggio del mio plotone che mi ordinava di fare rapporto al capitano
Andrews dall'altra parte della base. Andrews era al comando di una
compagnia in un altro battaglione della nostra brigata. Il messaggio non
diceva ``subito'', ma era implicito. Quando trovai la tenda che Andrews
usava come quartier generale della compagnia, entrai e salutai:
«Specialista Bishop a rapporto come da ordini, signore».

«Bishop», il capitano alzò a stento gli occhi dal portatile, «riposo.
Sono certo che sai che abbiamo lasciato la Terra senza essere al
completo, se la gente non riusciva ad arrivare in Ecuador entro la data
di partenza, restava fuori.»

Annuii. Il trasporto in Ecuador era stato, ironia della sorte, più
difficile del trasporto fino all'orbita e attraverso le stelle. A causa
della carenza di personale, avevamo sergenti al comando di squadre,
sergenti scelti nel ruolo di primi sergenti, tenenti che agivano come
capitani ecc. L'altro gruppo di fuoco nella mia squadra era composto
soltanto da tre soldati, e il battaglione aveva messo insieme nuove
squadre da unità sotto organico. Nel complesso, la X divisione aveva il
14\% in meno di personale autorizzato, senza neppure contare i
battaglioni di artiglieria da campo che erano stati deliberatamente
esclusi. O la maggior parte della brigata di supporto, che era stata a
sua volta esclusa, quindi non avevamo i nostri ingegneri di
combattimento con noi. Era un grosso problema che la divisione stava
ancora risolvendo.

«Ho qui il tuo fascicolo personale», continuò Andrews e io sbiancai.

Avevo visto il mio fascicolo personale, quello originale in cartaceo,
ancora nella triplice copia ufficiale dell'esercito. All'uscita dalla
Nigeria, si leggeva ``cattivo comportamento'', tutto maiuscolo e
sottolineato due volte. Non c'era nemmeno una faccina sorridente al
posto del punto sopra la I di ``cattivo'', quindi sapevo che era grave.
Speravo che il capitano Andrews non l'avrebbe sottolineato di nuovo e
cerchiato con una grande faccia arrabbiata. Nella mia corsa frenetica a
ritroso nella memoria, non riuscivo a pensare a niente che avessi
incasinato più del solito da quando avevamo lasciato la Terra. A
controbilanciare il mio presunto cattivo comportamento c'era un Cuore
viola al merito, che spiegava le cicatrici sul mio braccio sinistro, e
le scartoffie avevano cominciato a propendere per conferire una Stella
di bronzo. Le scartoffie si erano fermate al livello di brigata, cosa
che capivo considerando le circostanze. Comunque apprezzai il gesto.

«Conosci il sergente scelto Agnelli?», chiese Andrews, e io annuii di
nuovo. Poi disse qualcosa di inaspettato. «C'è un posto vacante per un
sergente nella sua squadra e tu sei stato raccomandato.» Non precisò chi
mi aveva raccomandato. «È tuo, se lo vuoi. Non abbiamo tempo, o risorse,
per un regolare programma di formazione, perciò imparerai sul campo. Se
t'interessa, fai rapporto ad Agnelli e io mi occuperò delle scartoffie.»

Merda. Non sapevo se ero pronto a comandare un gruppo di fuoco, tre tizi
che non conoscevo. Andrew si schiarì la gola: «Sono io che aspetto una
risposta, Bishop. So che è insolito, ma non abbiamo tempo per la normale
procedura qui».

«Ah, sì, signore. Grazie. Voglio dire, sì, sono onorato. Sì.» Una
promozione avrebbe dovuto essere un processo più formale, ero stato
preso del tutto alla sprovvista. La divisione aveva buchi che dovevano
essere tappati.

«Sì, sì.» Andrews non sembrava molto convinto. «Agnelli ti tirerà fuori
dai guai. E, Bishop?»

«Sì?»

«Stai lontano dai furgoni dei gelati.» Non sorrideva.

\hfill\break

Prendere il comando di un nuovo gruppo di fuoco non fu la parte più
difficile della mia promozione, i miei tre ragazzi avevano tutti
esperienza in Nigeria e sapevano cosa fare senza che io avessi bisogno
di dire loro molto. Il sergente scelto Agnelli era paziente con me e, al
di là della parte amministrativa del lavoro che all'inizio era
opprimente, non ebbi molti problemi a adattarmi. L'esercito non ebbe il
tempo di sottopormi all'addestramento standard, il che mi salvò dal
memorizzare un mucchio di stronzate, ed ero certo che mi si sarebbe
ritorto contro dopo. La parte più difficile della promozione fu lasciare
la mia vecchia unità, il Sergente Koch, Ski e soprattutto Cornpone. Lui
aveva una lacrima nell'occhio, che disse essere stata causata dalla
sabbia che si sollevava, e anche a me s'inumidirono gli occhi. «Cazzo»,
disse Cornpone, «pensa a quante cose stupide fai quando ci sono io. Cosa
farai senza di me? E l'esercito ti ha messo a capo di un gruppo di
fuoco?» Scosse la testa con tristezza.

«Se vedrò un furgone dei gelati correrò nella direzione opposta.»

Sbuffò: «Bish, questo risolve, tipo, il 10\% del problema. Ti trovi in
una situazione e dovresti pensare: ``Cosa farebbe Jesse?''».

«E poi fare l'opposto?» Risi.

Fece un'alzata di spalle esagerata, poi mi abbracciò, dandomi pacche
sulla schiena. «Abbi cura di te e non dimenticare di scrivere.»

Tirai fuori il mio zPhone. «Scrivere? Con questo coso posso vedere la
tua brutta faccia ogni volta che voglio.»

«Oh, cazzo», strabuzzò gli occhi, «l'avevo dimenticato.» Tirò fuori il
suo zPhone. «Come si fa a bloccare la gente?»

\hfill\break

I mezzi di comunicazione kristang non erano più veloci della luce, ma le
voci di corridoio dell'esercito degli Stati Uniti sì. «Sergente Bishop?»
Shauna fu la prima persona a chiamarmi in merito alla mia promozione.
«Ragazzi, promuoveranno proprio tutti.»

«Difatti.»

«Sul serio, buon per te, Bish.»

«Vuoi la verità? Non ho idea di cosa sto facendo.»

«Bish, combatteremo contro gli alieni. Nessuno sa cosa sta facendo. Ce
lo inventeremo sul campo.»

Cazzo, era bello sentirla. Il che mi ricordò che dovevo rivedere le
regole dell'esercito sulla fraternizzazione: Shauna non era sotto il mio
comando, era in un altro battaglione. Significava che andava bene se
continuavo a stare con lei? Di sicuro lo speravo.

\hfill\break

La mia nuova squadra era guidata dal sergente scelto Salvatore Agnelli
che, a dispetto del nome, aveva i capelli biondi e sembrava più tedesco
che italiano. Col mio gruppo di fuoco avevo avuto fortuna, e sia io che
Agnelli lo sapevamo. I soldati Chen e Baker e lo specialista Sanchez
erano tutti veterani della Nigeria, e Sanchez aveva prestato servizio
per un po' in Corea prima di allora. Greg Chen era un Abc dalla
periferia del Maryland, e dovette spiegare che Abc significava American
Born Chinese, cinese nato in America, dal momento che non avevo idea di
cosa intendesse. Jeron Baker era cresciuto in una fattoria in Alabama,
di cui era immensamente orgoglioso, e dalla quale era immensamente grato
di allontanarsi. E Pete Sanchez veniva dal Kansas orientale, i genitori
avevano una piccola fattoria, ma la madre lavorava come infermiera e il
padre in una concessionaria di auto, la loro fattoria non era abbastanza
grande per provvedere al sostentamento di una famiglia. Il loro
precedente sergente era adesso sergente scelto, a capo di una squadra
della II brigata. Iniziammo alla grande, perché nessuno di loro fece
battute su Barney, né chiese qualcosa in proposito, né sembrò in qualche
modo incuriosito da quello che avevo fatto la prima volta che avevo
visto un ruhar. Avevano letto tutto, era storia vecchia e non
significava più nulla. Noi quattro ci conoscemmo durante gli
addestramenti e le esercitazioni al combattimento la settimana
successiva. Ogni notte collassavo sul letto del tutto esausto e felice
come non lo ero da tempo. Avevo una grande squadra e ci stavamo
perfezionando per la battaglia contro un nemico che aveva minacciato il
nostro pianeta. Quello che ci serviva era una vera missione.

Otto giorni dopo la mia promozione a sergente, la nostra brigata
ricevette l'ordine di schierarsi. Mi stavo sforzando di mettermi in pari
per la successiva simulazione di guerra, l'operazione Rasoio, che ci
sarebbe stata di lì a dieci giorni. Ricevemmo informazioni dal capitano
Teller, che aveva riunito la compagnia in un hangar vuoto ai confini
della base. Stava in piedi su un palco in fondo, con noi che ci
accalcavamo per ascoltarlo. «Gente, abbiamo un'opportunità.»

Chen gemette e sussurrò: «Bohica¹¹, amico». Era un detto popolare
nell'esercito, soprattutto quando un ufficiale diceva che avevamo
``un'opportunità''. Di solito significava che il peggio doveva ancora
venire e che saremmo stati fregati in un modo o nell'altro, di qui
l'espressione radicata nel tempo.

Teller aspettò un momento perché l'inevitabile mormorio morisse da sé.
«Abbiamo l'opportunità di mettere in azione i nostri nuovi giocattoli e
fare qualcosa di utile in questa guerra. L'Unef ha ricevuto ordini,
partiamo tra due giorni e un risveglio. La I brigata entrerà con un
battaglione inglese e uno francese. L'Unef è tutta coinvolta, saremo al
completo. L'intera forza leverà le tende nelle prossime due settimane.»

Un tizio in prima fila prese la parola: «Qual è la missione, capitano?
Prenderemo a calci in culo qualche roditore?». La domanda fu posta con
serietà mortale.

Teller fece la faccia di chi ha appena morso un limone: «No, i kristang
dicono che non siamo pronti per questo. Ho visto simulazioni di
battaglia dei kristang, e hanno ragione, non siamo pronti. Saremmo solo
d'impaccio. La maggior parte dei combattimenti avviene in orbita o più
lontano e, dato che non abbiamo una Marina spaziale, non possiamo fare
nulla. I kristang hanno riconquistato un mondo coloniale che avevano
perso contro i ruhar tempo fa e li stanno cacciando dal pianeta. Hanno
raggiunto una tregua, l'accordo permette ai ruhar di evacuare la
popolazione civile in un periodo di tredici mesi, ci vorrà così tanto
perché ci sono quasi un milione di ruhar ora. Il nostro lavoro è
occupare il pianeta e facilitare l'evacuazione».

Questa notizia non fece piacere alla compagnia.

«Servizio di guarnigione?!»

«Merda. Faremo da corpi di pace armati.»

«Dovremo indossare dei dannati caschi blu?»

Teller tenne le mani alzate finché non ci calmammo. «Sono venuto qui per
combattere i ruhar, proprio come voi, per uccidere quei figli di puttana
con la faccia da topo. Non posso fare quello che voglio, sono
nell'esercito. Il nostro lavoro qui è difendere la Terra e, visto che
abbiamo dimostrato di non poterlo fare da soli, il modo migliore per
proteggere la gente a casa è lavorare con i nostri alleati. Le forze
kristang sono piuttosto assottigliate in questo momento, non hanno molte
truppe per mantenere i siti che hanno conquistato. Ecco perché hanno
chiesto la fanteria. C'è una specie di pericolo biologico su questo
pianeta, che impedisce ai kristang di atterrare subito; dicono che i
ruhar stanno tradendo la Regola Numero Quattro. I ruhar dicono che il
virus, o qualunque cosa sia, risale a quando i kristang avevano occupato
il pianeta. Di questo sta a loro discutere. Il punto è che i kristang
non possono far atterrare le loro truppe su questo pianeta, a meno che
non indossino tute di protezione complete, e non succederà. Fino a
quando i kristang non saranno in grado di sviluppare un antidoto al
virus e distribuirlo in quantità, l'Unef sarà responsabile in pieno a
terra. Questo è un lavoro che possiamo fare, e lo faremo al meglio delle
nostre capacità. Questi sono i nostri ordini.»

«Quali sono le regole d'ingaggio, signore?», chiese qualcuno ad alta
voce dalle ultime file.

«Non sparare, a meno che non siano loro a sparare per primi. Che ci
piacciano o no, i ruhar su questo pianeta sono per lo più agricoltori,
il che significa che tra loro ci sono donne e bambini. I kristang dicono
di non aspettarci molta opposizione, i civili ruhar dovrebbero essere
felici di essere vivi e di poter tornare a casa.» Ci fu un sacco di
brontolio dalla compagnia Able. Teller non fece una piega. «Questo è un
test, nel caso non l'abbiate ancora capito. Finora, tutto ciò che noi
umani abbiamo mostrato ai kristang è la nostra morte, senza infliggere
molti danni ai ruhar.» Si fermò, si accigliò e continuò: «Durante la
seconda guerra mondiale», pronunciò le parole con precisione,
«l'esercito americano non colpì le spiagge della Normandia nella prima
battaglia. Atterrammo prima in Nord Africa, e quelli di voi che
conoscono la loro storia sanno che è una cosa grandiosa. Non eravamo per
niente in grado di affrontare i tedeschi nel 1942, se fossimo atterrati
prima in Normandia, ci saremmo fatti prendere a calci in culo per tutta
la Manica. Combattendo in Nord Africa, in Sicilia e in Italia, imparammo
a contrastare i tedeschi, scoprimmo che cosa funzionava. I kristang ci
hanno dato alcuni nuovi giocattoli scintillanti con cui combattere,
questo non significa che conosciamo le migliori tattiche per usarli
contro i ruhar in combattimento a terra. I kristang ci hanno affidato
quello che dovrebbe essere un lavoro semplice, non possiamo rovinare
tutto. Prendiamo il controllo del pianeta, spostiamo i civili ruhar
verso gli ascensori spaziali e portiamo i loro sederi pelosi fuori dal
pianeta, mettiamo in sicurezza il posto, finché i kristang non
risolveranno le cose di sopra e potranno far atterrare le loro truppe.
Mostreremo ai kristang che le Forze di spedizione delle Nazioni unite
sono disciplinate, competenti e pronte ad affrontare le operazioni di
combattimento in futuro. Capi plotone, voglio le vostre squadre pronte
per l'ispezione di partenza alle 14:00 di dopodomani. Ci sono altre
domande?».

Presi la parola: «Signore, dove stiamo andando e che condizioni possiamo
aspettarci?». Volevo sapere il più possibile rispetto a quello in cui si
sarebbe cacciato il mio gruppo di fuoco.

«La risposta breve è che stiamo andando in culo a Nettuno.» Il sorriso
gli sollevò solo un lato della bocca. «Atmosfera ossigeno-azoto, un po'
più ossigeno di quello cui siamo abituati. La gravità è maggiore del 2\%
a quella sulla Terra. La stella è più calda del nostro Sole, ma il
pianeta è più lontano, quindi il clima è simile a quello della Terra:
caldo all'equatore, calotte glaciali ai poli. È per lo più una colonia
agricola, i ruhar chiamano il posto Gehtanu, che si traduce più o meno
come ``Nuovo campo di cereali'', o qualcosa del genere.» Questo scatenò
una risatina dalla folla. «I kristang lo chiamano Pradassis. Noi lo
chiameremo Paradiso. Ci sono un continente enorme, un paio di
arcipelaghi e grandi isole, come la Groenlandia, il resto è oceano. I
nostri PowerPoint ranger hanno messo insieme un pacchetto informativo
che i vostri capi plotone dovrebbero avere su tablet da distribuire, vi
si legge che il continente principale è in gran parte coperto da
praterie piane, come le Grandi Pianure degli Stati Uniti o le steppe
russe. Molte vaste fattorie sono in parte gestite da robot, e i ruhar
sono raggruppati in villaggi. Tutte le informazioni che abbiamo saranno
disponibili sui vostri notepad a breve, ho bisogno che tutti studino
durante il volo. L'intera X divisione, più i nostri amici inglesi e
francesi, s'imbarcherà su tre navi da trasporto e il viaggio durerà
sedici giorni. Dopo il lancio, i kristang terranno un cacciatorpediniere
più due fregate in orbita per il supporto al fuoco, ma per oggi è tutto.
Il generale Meers ha assicurato ai kristang che possiamo farcela, noi
non vogliamo dimostrargli che si sbaglia.»

\hfill\break

Shauna mi chiamò prima che potessi contattarla io. «Hai sentito? Ce ne
andiamo!»

«Sì, andiamo su Paradiso.»

«Paradiso!» Rise. Non stavamo insieme dalla mia promozione a sergente,
eravamo entrambi troppo impegnati, e io troppo stanco. Mi venne in mente
che avrei potuto non rivederla mai più: era probabile che, una volta
lasciata la nostra base su Campo Alfa, saremmo stati su diverse navi da
trasporto, poi assegnate a diversi siti sparsi sulla superficie di
Paradiso.

«Ehi, ehm, buona fortuna con la qualifica per la fanteria.» Non volevo
chiudere, ma non sapevo cosa dire.

«Grazie, dovrà aspettare finché non arriviamo là.» C'era gente che
parlava ad alta voce in sottofondo. «Joe, devo andare, restiamo in
contatto?»

«Certo.» E questo è tutto. Almeno saremmo stati sullo stesso pianeta.

Potreste pensare che ``due giorni e un risveglio'' significhi che
trascorsero due giorni interi, più una fastosa notte di sonno pieno,
prima che ci dovessimo radunare per la partenza. Vi sbagliate di grosso.
Il ``risveglio'' avvenne a un'ora improbabile del mattino, avevo
l'impressione che la mia testa avesse a stento toccato il cuscino prima
che le luci scattassero nella tenda e i miei piedi rimbalzassero sul
terreno. Alzarmi al canto del gallo non era una novità per me e non mi
dispiaceva, anche se non ero per natura una persona mattiniera. E
nemmeno una persona che si alzava a metà mattina. Mia sorella era una
persona mattiniera fastidiosamente allegra e, crescendo, più di una
volta avrei voluto soffocarla con un cuscino.

Comunque, dovetti svegliarmi a un'ora antelucana, il che significa che
poi dovetti svegliare il mio gruppo di fuoco. La differenza tra alzarsi
a un'ora antelucana e alzarsi al canto del gallo è che alzarsi a un'ora
antelucana significa che ti alzi dal letto, prepari tutta l'attrezzatura
e l'assembli per l'ispezione di partenza solo per startene in piedi tre
cazzo di ore, quando avresti potuto usare quelle tre ore per qualcosa di
produttivo, tipo dormire. Da sergente, adesso dovevo sopportare gli
sguardi risentiti della mia squadra, mantenendo un'espressione neutrale
sul viso.

Tutta la partenza fu un casino. I kristang volevano che ci imbarcassimo
sulle navicelle nel modo più rapido ed efficiente possibile, il che
comportava ammassare il numero massimo di persone a bordo, in una corsa
frenetica, anche se ciò significava dividere plotoni, squadre o persino
gruppi di fuoco. Quando alla fine salimmo a bordo di una navicella,
Baker e Chen avevano i piedi nella porta quando l'addetta della polizia
militare stese una mano per impedire a chiunque altro di salire. «Siamo
al completo, passiamo alla prossima nave», annunciò.

Tirai fuori con gesto frenetico Baker e Chen, afferrandoli per lo zaino
e feci cenno ad altri due ragazzi di venire avanti. Cazzo, non avrei mai
permesso che il mio nuovo gruppo di fuoco si dividesse subito dopo che
ci eravamo conosciuti. Penserete che saremmo stati i primi della fila a
salire a bordo della navicella successiva, ma gli agenti che si
occupavano dell'organizzazione avevano altro in mente e diverse
navicelle decollarono prima che ci facessero cenno di andare avanti.
Essere a bordo di una navicella kristang non era poi l'ultimo dei
casini: quando la nostra attraccò e ci imbarcammo, non solo scoprimmo
che eravamo a bordo di una nave diversa dal resto del nostro plotone, ma
anche che era per la maggior parte occupata dalla III divisione di
fanteria dell'esercito americano, non dalla nostra unità di
appartenenza, il X. La nave conteneva persone del III, del X, la maggior
parte di una compagnia di Marine degli Stati Uniti, più un'intera
compagnia di cinesi, due squadre di inglesi e una manciata di soldati
indiani. O i kristang non capivano il concetto di carico di dislocamento
da combattimento, oppure avevano così poco rispetto per le capacità di
combattimento umane da pensare che non avesse importanza che le nostre
unità fossero sparpagliate per mari e per monti. È probabile che il
quartier generale dell'Unef stesse tenendo traccia di chi fosse dove, ma
non avevo modo di contattare nessuno nella mia catena di comando.
Ottenni che dormissimo con qualcun altro del X. Fu un viaggio solitario
e fui riconoscente che i miei ragazzi l'avessero presa in modo sportivo,
con una scrollata di spalle. Se quello fosse stato il peggior casino in
cui ci saremmo cacciati, saremmo stati benissimo.